\subsection{等变函数}

在本小节中，我们取定 G 的一个闭子群 H。用 \(\varpi\) 表示 G 到 G/H 上的典范投射。

除 G 上的哈尔测度 \(\mu\) 外，我们再取定 H 上的一个左哈尔测度 \(\beta\)。

根据 INT, VII, 第 56 页, § 2, 第 5 目, 定理 2，存在连续函数 \(\kappa\colon\mathbf{G}\to\mathbf{R}_{+}^{*}\)，使得 \(\kappa(xh)=\Delta_{\mathrm{H}}(h)\Delta_{\mathrm{G}}(h)^{-1}\kappa(x)\) 对每个 \((x,h)\in\mathbf{G}\times\mathbf{H}\) 成立。我们取定这样一个函数 \(\kappa\)。用 \(\nu\) 表示 G/H 上的测度 \((\kappa\cdot\mu)/\beta\)；由出处同上，它是 G 下拟不变的非零正测度。它的支撑等于 G/H（INT,

VII, 第 10 页, § 1, 第 1 目）。根据 INT, VII, 第 43 页, § 2, 第 2 目, 命题 4，测度 \(\nu\) 是 G/H 上唯一的使得 G 上的测度 \(\nu^{\sharp}\) 等于 \(\kappa \cdot \mu\) 的测度。我们把空间 G/H 赋以测度 \(\nu\)（因而例如我们写 \(\mathcal{L}^{p}(\mathrm{G / H}) = \mathcal{L}^{p}(\mathrm{G / H},\nu)\)）。

 我们称 G 的子集 \(S \subset G\) 是模 H 零测的，如果 \(\varpi(S)\) 是 \(\nu\) -零测的。这一条件不依赖于 G 与 H 上哈尔测度的选取。它蕴含 S 是局部 \(\mu\) -零测的（INT, VII, 第 47 页, § 2, 第 3 目, 命题 6, a)）。我们称 G 的元素所具有的一个性质模 H 几乎处处为真，如果 G 中使该性质不成立的元素所成的集是模 H 零测的。

设 \(\pi\) 是 H 在希尔伯特空间 E 中的一个酉表示。用 \(\mathcal{F}_{\pi}(\mathrm{G})\) 表示 G 上取值于 E 的函数 \(f\) 所成的向量空间，其中 \(f(xh) = \pi(h)f(x)\) 对一切 \((x,h) \in \mathrm{G} \times \mathrm{H}\) 成立。我们称 \(\mathcal{F}_{\pi}(\mathrm{G})\) 的元素为 G 上的 \(\pi\) -等变函数。

G 在 C 上的平凡表示所对应的空间 \(\mathcal{F}_{1}(\mathrm{G})\) 与 G/H 上复值函数的空间 \(\mathcal{F}(\mathrm{G / H})\)，借助映射 \(f\mapsto f\circ \varpi\)（即 \(\mathcal{F}(\mathrm{G / H})\) 到 \(\mathcal{F}_1(\mathrm{G})\) 中的映射）等同。

对任意函数 \(f\)（属于 \(\mathcal{F}_{\pi}(\mathrm{G})\)），函数 \(\|f\|\) 属于 \(\mathcal{F}_{1}(\mathrm{G})\)，因为 \(\pi\) 是酉的。这使我们可以把 \(\|f\|\) 看作 \(\mathrm{G/H}\) 上的一个函数，例如我们将写作 \(\|f(x\mathrm{H})\|\)，以表示该函数在元素 \(x\mathrm{H}\)（\(\mathrm{G/H}\) 的元素）处的值。

设 \(f\) 是 \(\mathcal{F}_{\pi}(\mathrm{G})\) 中的一个函数；若 \(\|f\|\) 在 G/H 的一个紧子集外为零，则称它模 H 在一个紧集外为零。这等价于说，存在 G 的紧子集 K，使 \(f\) 在 \(\mathrm{K} \cdot \mathrm{H}\) 之外为零（TG, III, 第 33 页, 命题 10）。

用 \(\mathcal{K}_{\pi}(\mathrm{G})\) 表示 G 上属于 \(\mathcal{F}_{\pi}(\mathrm{G})\) 且模 H 具有紧支撑的连续函数所成的空间。

与 \(\mathcal{H}_{\pi}(\mathrm{G})\) 类似的一个空间，当 \(\pi\) 是 H 到 \(\mathbf{R}_{+}^{*}\) 中的连续同态时，已在 INT, VII, 第 39 页, § 2, 第 1 目 中定义过。

设 \(p \in [1, +\infty[\)。对每个 \(f \in \mathcal{F}_{\pi}(\mathrm{G})\)，记

\[
\mathrm{N} _ {p} (f) = \left(\int_ {\mathrm{G/H}} ^ {*} \| f \| ^ {p} d \nu\right) ^ {1 / p}.
\]

它是一个实数或 \(+\infty\)。用 \(\mathcal{F}_{\pi}^{p}(\mathrm{G},\nu)\)，或简作 \(\mathcal{F}_{\pi}^{p}(\mathrm{G})\)，表示由函数 \(f\in \mathcal{F}_{\pi}(\mathrm{G})\) 中使 \(\mathrm{N}_p(f)\)

为有限者所成的子空间。空间 \(\mathcal{F}_{\pi}^{p}(G)\) 赋以映射 \(N_{p}\) 后是一个半赋范空间。

空间 \(\mathcal{K}_{\pi}(\mathrm{G})\) 包含于 \(\mathcal{F}_{\pi}^{p}(\mathrm{G})\)。它在 \(\mathcal{F}_{\pi}^{p}(\mathrm{G})\) 中的闭包记作 \(\mathcal{L}_{\pi}^{p}(\mathrm{G},\nu)\)，或简作 \(\mathcal{L}_{\pi}^{p}(\mathrm{G})\)；我们称它为 G 上 \(\pi\) -等变函数中模 H 以指数 \(p^{e}\) 可积的那些函数所成的空间。

下面命题的各断言，当 \(\pi\) 是一维平凡表示时，是 INT, IV, 第 128 页, § 3, 第 3 目, 命题 6 及第 131 页, § 3, 第 4 目, 定理 3 的推论。

命题 3. --- a) 设 \((f_n)_{n \in \mathbf{N}}\) 是 \(\mathcal{F}_\pi^p(\mathrm{G})\) 中的一个序列，使得级数

\[
\sum_ {n = 0} ^ {+ \infty} \mathrm{N} _ {p} (f _ {n})
\]

收敛。于是以 \(f_n(g)\) 为通项的级数对位于一个模 H 零测集 T 之外的每个 \(g\) 绝对收敛。设 \(f\) 是 G 上取值于 E 的函数，它在 G−T 上等于该级数的和，而在 T 上等于 0。这时我们有 \(f \in \mathcal{F}_\pi^p(\mathrm{G})\)，且以 \(f_n\) 为通项的级数收敛于 \(f\)，这里收敛是在空间 \(\mathcal{F}_\pi^p(\mathrm{G})\) 中取的；

\begin{enumerate}
\def\labelenumi{\alph{enumi})}
\setcounter{enumi}{1}
\item
  设 \((f_n)\) 是 \(\mathcal{L}_{\pi}^{p}(\mathrm{G})\) 中收敛于一个函数 \(f\) 的序列。存在 \((f_{n_k})\)——它是 \((f_n)\) 的一个子列——使得 \(f_{n_k}(g)\) 收敛于 \(f(g)\)，对位于一个模 H 零测集之外的每个 \(g\) 成立；
\item
  半赋范空间 \(\mathcal{F}_{\pi}^{p}(\mathrm{G})\) 与 \(\mathcal{L}_{\pi}^{p}(\mathrm{G})\) 是完备的。
\end{enumerate}

在断言 \(a\) ) 中，说以 \(f_n\) 为通项的级数收敛于 \(f\)——这里是指在空间 \(\mathcal{F}_{\pi}^{p}(\mathrm{G})\) 中——意思是部分和的序列 \(f_0 + \cdots + f_n\) 收敛于 \(f\)，收敛是在 \(\mathcal{F}_{\pi}^{p}(\mathrm{G})\) 中取的。这时也称 \(f\) 是该级数的一个和。

我们来证明 \(a\) )。根据 INT, IV, 第 128 页, § 3, 第 3 目, 命题 6，存在一个 \(\nu\) -零测集 S ⊂ G/H，使得以 \(\|f_n(gH)\|\) 为通项的级数对 gH \(\notin\) . S 绝对收敛。此外，对于 gH \(\notin\) . S 等于该级数之和、对于 gH \(\in\) . S 为零的函数 h 满足 N \(_p(h)\) \textless{} + \(\infty\)。

集合 T = \(\overline{\varpi}(S)\) 是模 H 零测的。对每个 \(g \notin T\)，以 \(f_{n}(g)\) 为通项的级数在 E 中绝对收敛。令 \(f(g) = \sum f_{n}(g)\)（对 \(g \notin T\)），否则令 \(f(g) = 0\)。我们有 \(f \in \mathcal{F}_{\pi}(G)\)。注意 \(\|f(gH)\| \leqslant h(gH)\) 对所有 \(g \in G\) 成立，由此 \(N_{p}(f) \leqslant\)

\(N_{p}(h) < +\infty\)。于是我们有 \(f \in \mathcal{F}_{\pi}^{p}(G)\)。类似地，由此得出

\[
\mathrm{N} _ {p} \Big (f - \sum_ {n = 0} ^ {k} f _ {n} \Big) \leqslant \sum_ {n = k + 1} ^ {+ \infty} \mathrm{N} _ {p} (f _ {n})
\]

对所有 \(k \in \mathbf{N}\) 成立，因此级数 \(\sum f_n\) 收敛于 \(f\)，收敛是在 \(\mathcal{F}_{\pi}^{p}(\mathrm{G})\) 中取的。断言 \(a)\) 得证。

我们来证明 \(b\) )。序列 \((\|f_n - f\|)_{n \in \mathbf{N}}\) 在空间 \(\mathcal{L}^p(\mathrm{G}/\mathrm{H})\) 中收敛于 0。根据 INT, IV, 第 131 页, § 4, 第 3 目, 定理 3，存在一个子序列 \((\|f_{n_k} - f\|)_{k \in \mathbf{N}}\)，它 \(\nu\) -几乎处处收敛于 0。于是序列 \((f_{n_k}(g))_{k \in \mathbf{N}}\) 收敛于 \(f(g)\)，对位于一个模 H 零测集之外的每个 \(g\) 成立。

最后证明 \(c\) )。设 \((f_n)_{n \in \mathbf{N}}\) 是 \(\mathcal{F}_\pi^p(\mathrm{G})\) 中的一个柯西序列。存在严格递增的整数序列 \((n_k)_{k \in \mathbf{N}}\)，使得 \(\mathrm{N}_p(f_{n_{k+1}} - f_{n_k}) \leqslant 2^{-k}\) 对所有 \(k \in \mathbf{N}\) 成立。对每个 \(k \in \mathbf{N}\)，令 \(h_k = f_{n_{k+1}} - f_{n_k} \in \mathcal{F}_\pi^p(\mathrm{G})\)。由 \(a\) )，以 \(h_k\) 为通项的级数在 \(\mathcal{F}_\pi^p(\mathrm{G})\) 中收敛；记其和为 \(h\)。对每个 \(\ell \in \mathbf{N}\)，由此得出

\[
f _ {n _ {0}} + \sum_ {k = 0} ^ {\ell} h _ {k} = f _ {n _ {\ell + 1}},
\]

于是 \(f_{n_{0}} + h\) 是序列 \((f_{n})\) 的一个聚点。从而该序列收敛（TG, II, 第 14 页, 命题 5 的推论 2）。这样，空间 \(\mathcal{F}_{\pi}^{p}(\mathrm{G})\) 是完备的；又因空间 \(\mathcal{L}_{\pi}^{p}(\mathrm{G})\) 在 \(\mathcal{F}_{\pi}^{p}(\mathrm{G})\) 中是闭的，它也是完备的。

推论. --- 设 \((f_{n})_{n\in\mathbf{N}}\) 是 \(\mathcal{L}_{\pi}^{p}(\mathrm{G})\) 中的一个柯西序列，并设 \(f\in\mathcal{F}_{\pi}(\mathrm{G})\) 使得 \(f_{n}(g)\) 模 H 几乎处处收敛于 \(f(g)\)。则 \(f\in\mathcal{L}_{\pi}^{p}(\mathrm{G})\)，且 \((f_{n})_{n\in\mathbf{N}}\) 在 \(\mathcal{L}_{\pi}^{p}(\mathrm{G})\) 中收敛于 f。

事实上，函数 \(f\) 与序列 \((f_n)\) 在 \(\mathcal{L}_{\pi}^{p}(\mathrm{G})\) 中的极限模 H 几乎处处重合。

用 \(\mathrm{L}_{\pi}^{p}(\mathrm{G})\) 表示与半赋范空间 \(\mathcal{L}_{\pi}^{p}(\mathrm{G})\) 相伴随的分离赋范拓扑向量空间；它是一个巴拿赫空间。

设 \(f\) 是 \(\mathrm{G}\) 上取值于 \(\mathrm{E}\) 的一个函数。用 \(\mathrm{S}_{f}\) 表示由 \(x\in \mathrm{G}\) 中使函数 \(h\mapsto f(xh)\)（定义在 \(\mathrm{H}\) 上）不属于 \(\mathcal{L}_{\mathrm{E}}^{1}(\mathrm{H})\) 者所成的集。我们有 \(\mathrm{S}_f\cdot h = \mathrm{S}_f\) 对所有 \(h\in \mathrm{H}\) 成立。

设 \(x \in \mathrm{G} - \mathrm{S}_f\)。由于 \(\pi\) 是酉表示，函数 \(h \mapsto \pi(h)^* f(xh)\) 可测（它是 H 上的映射 \(h \mapsto (h, f(xh))\)——后者到 \(\mathrm{H} \times \mathrm{E}\) 中可测——与 H × E 到 E 中的连续映射 \((g,x)\mapsto\pi(g)^{*}x\) 的复合，参见 INT, IV, 第 174 页, § 5, 第 3 目, 定理 1），并且它在 H 上可积。

我们通过下式在 G 上定义一个函数 \(f^{\pi}\)：

\[
f ^ {\pi} (x) = \int_ {\mathrm{H}} \pi (h) ^ {*} f (x h) d \beta (h)\tag{3}
\]

其中 \(x \in \mathrm{G} - \mathrm{S}_f\)；而 \(f^\pi(x) = 0\)（当 \(x \in \mathrm{S}_f\) 时）。

命题 4. --- 设 \(f \in \mathcal{L}_{\mathrm{E}}^{1}(\mathrm{G})\)。

\begin{enumerate}
\def\labelenumi{\alph{enumi})}
\item
  集合 \(\mathrm{S}_f\) 模 H 零测，且 \(f^\pi \in \mathcal{F}_\pi(\mathrm{G})\)；
\item
  设 C 是 G 的一个紧子集。若 f 连续且支撑含于 C，则 \(S_{f}\) 为空，函数 \(f^{\pi}\) 属于 \(\mathcal{H}_{\pi}(\mathrm{G})\) 且其支撑含于 \(C \cdot H\)。
\end{enumerate}

\(a\) ) 的第一部分由 INT, VII, 第 57 页, § 2, 第 5 目, c) 得出。设 \(w \in H\)。若 \(x \in G - S_f\)，则 \(xw \in G - S_f\)，且

\[
\begin{array}{c} f ^ {\pi} (x w) = \int_ {\mathrm{H}} \pi (h) ^ {*} f (x w h) d \beta (h) \\ = \int_ {\mathrm{H}} \pi (w ^ {- 1} y) ^ {*} f (x y) d \beta (y) = \pi (w) f ^ {\pi} (x), \end{array}
\]

而若 \(x \in S_{f}\)，则 \(f^{\pi}(xw) = 0 = \pi(w)^{*}f^{\pi}(x)\)。因此函数 \(f^{\pi}\) 属于 \(\mathcal{F}_{\pi}(\mathrm{G})\)。

现在设 \(f\) 连续且支撑含于 C。于是对每个 \(x \in G\)，映射 \(h \mapsto f(xh)\) 属于 \(\mathcal{K}(H)\)，从而可积，这就证明了 \(S_f = \varnothing\)。

我们来证明 \(f^{\pi}\) 连续。设 \(x \in G\)。设 U 是 \(x\) 在 G 中一个相对紧的开邻域。

对每个 \(y \in \mathrm{U}\)，我们有

\[
\begin{array}{l} \| f ^ {\pi} (y) - f ^ {\pi} (x) \| \leqslant \int_ {\mathrm{H}} \| f (y h) - f (x h) \|   d \beta (h) \\ \qquad = \int_ {\mathrm{H} \cap (y ^ {- 1} \mathrm{C} \cup x ^ {- 1} \mathrm{C})} \| f (y h) - f (x h) \|   d \beta (h) \\ \qquad \leqslant \beta (\mathrm{H} \cap \mathrm{U} ^ {- 1} \mathrm{C}) \sup _ {h \in \mathrm{U} ^ {- 1} \mathrm{C}} \| f (y h) - f (x h) \|, \end{array}
\]

而 \(f^{\pi}\) 的连续性则由 f 在 G 上的一致连续性得出。

最后，若 \(x \in G\) 满足 \(f^{\pi}(x) \neq 0\)，则存在 \(h \in H\) 使 \(f(xh) \neq 0\)，于是 xh 属于 C，而 x 属于 C·H。由于 C·H 在 G 中闭，我们断定 \(f^{\pi}\) 的支撑含于 C·H。

 设 C 是 G 的一个紧子集。设 \(u \in \mathcal{K}_{+}(G)\) 是满足 \(u(x) > 0\)（对一切 \(x \in C\)）的函数。设 v 是 G 上由下式定义的函数

\[
v (x) = \int_ {\mathrm{H}} u (x h) d \beta (h)\tag{4}
\]

对所有 \(x \in G\) 成立；沿用前面的记号，我们有 \(v = u^{1}\)，它对应于 H 的一维平凡表示。函数 v 连续且为正；它属于 \(\mathcal{F}_{1}(G)\)，其支撑含于 \(\operatorname{Supp}(u) \cdot H\)，并且我们有

\[
\inf _ {x \in \mathrm{C} \cdot \mathrm{H}} v (x) > 0\tag{5}
\]

（INT, VII, 第 39--40 页, § 2, 第 1 目, 命题 1 及引理 1, a)）。

引理 5. --- 设 \(f \in \mathcal{F}_{\pi}(\mathrm{G})\) 是一个 \(\mu\) -可测且在 \(\mathrm{C} \cdot \mathrm{H}\) 之外为零的函数，使得函数 \(\|f\|\) \(\nu\) -可积（在 \(\mathrm{G}/\mathrm{H}\) 上）。设 s 是 \(\mathrm{G}\) 上取值于 E 的函数，满足

\[
 s (x) = \left\{ \begin{array}{l l} v (x) ^ {- 1} f (x) & \text{若 } x \in \mathrm{C} \cdot \mathrm{H} \\ 0 & \text{若 } x \in \mathrm{G} - \mathrm{C} \cdot \mathrm{H}. \end{array} \right.
\]

\begin{enumerate}
\def\labelenumi{\alph{enumi})}
\item
  函数 s 是 \(\mu\) -可测的；它属于 \(\mathcal{F}_{\pi}(\mathrm{G})\) 且在 C·H 之外为零；
\item
  函数 us 属于 \(\mathcal{L}^1 (\mathrm{G})\)，且模 H 几乎处处有 \((us)^{\pi} = f\)。
\end{enumerate}

函数 \(s\) 按定义在 C \(\cdot\) . H 之外为零；它是 \(\mu\) -可测的（INT, IV, 第 193 页, § 5, 第 10 目, 命题 16），因为函数 \(f\) \(\mu\) -可测且 \(v(x) > 0\) 对一切 \(x \in \mathrm{C} \cdot \mathrm{H}\) 成立。我们有 \(s \in \mathcal{F}_{\pi}(\mathrm{G})\)，因为 \(v \in \mathcal{F}_1(\mathrm{G})\)。

函数 \(f\) 是局部 \(\mu\) -可积的，因为 \(\|f\|\) 是 \(\nu\) -可积的函数（INT, VII, 第 47 页, § 2, 第 3 目, 命题 6, c)，注意测度 \(\nu^{\sharp} = \kappa \cdot \mu\) 等价于 \(\mu\)），从而由公式 (5)，函数 \(s\) 亦然。函数 \(us\) 可测且在 G 的一个紧子集之外为零；由于 \(u\) 有界，由此得出 \(us\) 是 \(\mu\) -可积的。

对每个 \(x\in \mathrm{G - S}_{us}\)，我们有

\[
\begin{array}{c} (u s) ^ {\pi} (x) = \int_ {\mathrm{H}} \pi (h) ^ {*} (u (x h) s (x h)) d \beta (h) \\ = \int_ {\mathrm{H}} u (x h) s (x) d \beta (h) = v (x) s (x), \end{array}
\]

因为 \(s(xh) = \pi(h)s(x)\)；最后一个断言由命题 4, a) 得出。

当 \(H = \{e\}\) 且 \(\pi\) 是一维时，下面的命题就是 INT, IV, 第 184 页, § 5, 第 6 目, 定理 5。

命题 5. --- 设 \(p \in [1, +\infty[\)。空间 \(\mathcal{L}_{\pi}^{p}(\mathrm{G})\) 是由满足如下条件的函数 \(f \in \mathcal{F}_{\pi}(\mathrm{G})\) 所成的空间：\(f\) \(\mu\) -可测，且函数 \(\|f\|\) 属于 \(\mathcal{L}^{p}(\mathrm{G}/\mathrm{H})\)。

设 \(f \in \mathcal{L}_{\pi}^{p}(G)\)。它是 \(\mathcal{L}_{\pi}^{p}(G)\) 中由 \(\mathcal{K}_{\pi}(G)\) 的元素组成的序列的极限。由命题 3, b) 及叶戈罗夫定理（INT, IV, 第 175 页, § 5, 第 4 目, 定理 2），函数 \(f\) 因此是 \(\mu\) -可测的；从而 G/H 上的函数 \(\|f\|\) 是 \(\nu\) -可测的（INT, VII, 第 47 页, § 2, 第 3 目, 命题 6, b)）。由于 \(N_{p}(f)\) 有限，函数 \(\|f\|\) 属于 \(\mathcal{L}^{p}(G/H)\)（INT, IV, 第 184 页, § 5, 第 6 目, 定理 5）。

 我们来证明逆命题。设 \(f\) 是一个 \(\mu\) -可测、属于 \(\mathcal{F}_{\pi}(\mathrm{G})\) 的函数，且满足 \(\|f\|\in\mathcal{L}^{p}(\mathrm{G}/\mathrm{H})\)。设 \(\varepsilon>0\)。我们来证明存在 \(\widetilde{f}\in\mathcal{K}_{\pi}(\mathrm{G})\)，使得

\[
\mathrm{N} _ {p} (f - \widetilde {f}) ^ {p} = \int_ {\mathrm{G/H}} ^ {*} \| f - \widetilde {f} \| ^ {p} d \nu \leqslant \varepsilon ,
\]

如此即完成证明。

 先设存在 G 的紧子集 C，使 \(f\) 在 \(\mathrm{C} \cdot \mathrm{H}\) 之外为零。设 \(u \in \mathcal{K}_{+}(\mathrm{G})\) 是满足 \(u(x) > 0\)（对一切 \(x \in \mathrm{C}\)）的函数，并按上式定义 \(v = u^{1}\)。用 \(\varphi\) 表示 \(u\) 的支撑的特征函数。

设 q 是 p 的共轭指数。若 p = 1，设 w 是 G 上等于 \(\sup_{x \in G} u(x)\) 的常值函数。若 p \textgreater{} 1，我们定义

\[
w (x) = \left(\int_ {\mathrm{H}} u (x h) ^ {q} d \beta (h)\right) ^ {1 / q}
\]

 对所有 \(x \in G\) 成立。在所有情形下，函数 w 都连续且为正；它属于 \(\mathcal{K}_{1}(G)\)（命题 4 应用于一维平凡表示及函数 \(u^{q}\)），从而在 G 上有界。我们令 \(W = \sup_{x \in G} w(x)\)。

 设 \(s\) 是 \(\mathbf{G}\) 上取值于 \(\mathbf{E}\) 的、把引理 5 应用于 \(f\) 时所得到的函数。令 \(g = \varphi s\)。函数 \(g\) 是 \(\mu\) -可测的，且满足 \(\| g \| \leqslant \| f \| / \inf_{x \in \mathrm{C} \cdot \mathrm{H}} v(x)\)。由于 \(g\) 在 \(u\) 的支撑之外为零且 \(\kappa\) 连续，我们有 \(g \in \mathcal{L}_{\mathrm{E}}^{p}(\mathrm{G}, \kappa \cdot \mu)\)。设 \(\widetilde{g}\) 是

\(\mathcal{K}_{\mathrm{E}}(\mathrm{G})\) 中的一个函数，使得

\[
\int_ {\mathrm{G}} ^ {*} \| g - \widetilde {g} \| ^ {p} \kappa d \mu \leqslant \frac {\varepsilon}{\mathrm{W} ^ {p}}.
\]

我们有 us = ug，从而模 H 几乎处处有 \(f = (us)^{\pi} = (ug)^{\pi}\)（引理 5, b)）。令 \(\widetilde{f} = (u\widetilde{g})^{\pi}\)；我们有 \(\widetilde{f} \in \mathcal{K}_{\pi}(\mathrm{G})\)（命题 4, b)）。对每个 \(x \in G - S_{ug}\)，我们得到

\[
\| (u g) ^ {\pi} (x) - (u \widetilde {g}) ^ {\pi} (x) \| \leqslant \int_ {\mathrm{H}} ^ {*} u (x h) \| g (x h) - \widetilde {g} (x h) \| d \beta (h)
\]

由此得

\[
\| (u g) ^ {\pi} (x) - (u \widetilde {g}) ^ {\pi} (x) \| ^ {p} \leqslant w (x) ^ {p} \int_ {\mathrm{H}} ^ {*} \| g (x h) - \widetilde {g} (x h) \| ^ {p} d \beta (h),
\]

这里在 \( p \) \textgreater{} 1 的情形用了赫尔德不等式。由于 \(S_{ug}\) 模 H 零测，由此得出

\[
\begin{array}{c} \int_ {\mathrm{G/H}} ^ {*} \| f - \widetilde {f} \| ^ {p} d \nu \leqslant \mathrm{W} ^ {p} \int_ {\mathrm{G/H}} ^ {*} \Big (\int_ {\mathrm{H}} ^ {*} \| g (x h) - \widetilde {g} (x h) \| d \beta (h) \Big) ^ {p} d \nu (x \mathrm{H}) \\ = \mathrm{W} ^ {p} \int_ {\mathrm{G}} ^ {*} \| g - \widetilde {g} \| ^ {p} \kappa d \mu \leqslant \varepsilon \end{array}
\]

这是根据 INT, VII, 第 46 页, § 2, 第 3 目, 命题 5, b)，它适用于 \(g - \widetilde{g}\) 在 G 的一个紧子集之外为零的情形。这就证明了当 \(f\) 模 H 在一个紧子集之外为零时所要求的性质。

现在考虑一般情形。由假设 \(\|f\|\in\mathcal{L}^{p}(G/H)\)，存在 G/H 的紧子集 L，使得

\[
\int_ {(\mathrm{G} / \mathrm{H}) - \mathrm{L}} ^ {*} \| f \| ^ {p} d \nu \leqslant \frac {\varepsilon}{2}
\]

（参见 INT, IV, 第 152 页, § 4, 第 6 目, 定理 4）。设 \(\varphi_{\mathrm{L}}\) 是 L 的特征函数，并令 \(f_{\mathrm{L}} = (\varphi_{\mathrm{L}} \circ \varpi)f\)。我们有

\[
\int_ {\mathrm{G} / \mathrm{H}} ^ {*} \| f - f _ {\mathrm{L}} \| ^ {p} d \nu = \int_ {(\mathrm{G} / \mathrm{H}) - \mathrm{L}} ^ {*} \| f \| ^ {p} d \nu \leqslant \frac {\varepsilon}{2}.
\]

函数 \(f_{\mathrm{L}}\) 是 \(\mu\) -可测的且模 H 在一个紧集之外为零。它属于 \(\mathcal{L}_{\pi}^{p}(\mathrm{G})\)，于是由前一种情形，存在 \(\widetilde{f} \in \mathcal{K}_{\pi}(\mathrm{G})\) 使得 \(\mathrm{N}_p(f_{\mathrm{L}} - \widetilde{f}) \leqslant (\frac{\varepsilon}{2})^{1/p}\)，从而

\[
\int_ {\mathrm{G/H}} ^ {*} \| f - \widetilde {f} \| ^ {p} d \nu \leqslant \varepsilon ,
\]

这正是所要证的。

 我们来考虑 \(p = 2\) 的情形。对 \(f_1\) 与 \(f_2\)（均属于 \(\mathcal{F}_{\pi}(\mathrm{G})\)），函数 \(x \mapsto \langle f_1(x) | f_2(x) \rangle\) 属于 \(\mathcal{F}_1(\mathrm{G})\)，因为表示 \(\pi\) 是酉的，从而通过取商定义了 \(\mathrm{G / H}\) 上的一个函数，我们和先前一样把它等同于 \(\langle f_1 | f_2 \rangle\)。我们有不等式 \(|\langle f_1 | f_2 \rangle| \leqslant \| f_1 \| \| f_2 \|\)（在 \(\mathcal{F}(\mathrm{G / H})\) 中）。

若 \(f_1\) 与 \(f_2\) 属于 \(\mathcal{L}_\pi^2(\mathrm{G})\)，则函数 \(\langle f_1 | f_2 \rangle\) 属于 \(\mathcal{L}_1^1(\mathrm{G})\)。特别地，它在 \(\mathrm{G/H}\) 上可积。映射

\[
(f _ {1}, f _ {2}) \mapsto \int_ {\mathrm{G} / \mathrm{H}} \langle f _ {1} | f _ {2} \rangle d \nu
\]

是 \(\mathcal{L}_{\pi}^{2}(\mathrm{G})\) 上的一个正埃尔米特形式；其相应的半范数是半范数 \(N_{2}\)。特别地，空间 \(\mathcal{L}_{\pi}^{2}(\mathrm{G})\) 是一个预希尔伯特空间，而 \(\mathrm{L}_{\pi}^{2}(\mathrm{G})\) 是与 \(\mathcal{L}_{\pi}^{2}(\mathrm{G})\) 相伴随的希尔伯特空间。

\subsection{诱导表示}

我们沿用前一日中关于 H 上的测度 \(\beta\) 与 G/H 上的测度 \(\nu\) 的记号和约定，以及函数 \(\kappa\colon\mathrm{G}\to\mathbf{R}_{+}^{*}\)。

存在一个连续函数 \(\eta\)，它定义在 \(\mathrm{G} \times \mathrm{G} / \mathrm{H}\) 上且取值于 \(\mathbf{R}_{+}^{*}\) 中，使得

\[
\eta (x, y \mathrm{H}) = \frac {\kappa (x y)}{\kappa (x)}
\]

对所有 \((x,y)\in\mathrm{G}\times\mathrm{G}\) 成立，且 \(\boldsymbol{\gamma}_{\mathrm{G}/\mathrm{H}}(x)\nu=(y\mapsto\eta(x^{-1},y))\cdot\nu\) 对 \(x\in G\) 成立（INT, VII, 第 56 页, § 2, 第 5 目, 定理 2, c)）。

设 \(p \in [1, +\infty[\)，并设 \(\pi\) 是 H 在复希尔伯特空间 E 中的一个酉表示。

引理 6. --- 设 \(f \in \mathcal{K}_{\overline{\pi}}(\mathrm{G})\)。对每个 \(g \in \mathrm{G}\)，函数

\[
\widetilde {f} \colon x \mapsto \eta (g ^ {- 1}, x \mathrm{H}) ^ {1 / p} f (g ^ {- 1} x)
\]

（它映 G 到 E 中）属于 \(\mathcal{K}_{\overline{\pi}}(\mathrm{G})\)，且满足 \(\mathrm{N}_p(\widetilde{f}) = \mathrm{N}_p(f)\)。

可以毫无困难地验证 \(\tilde{f} \in \mathcal{K}_{\overline{\pi}}(\mathrm{G})\)。由于

\[
\mathrm{N} _ {p} (\widetilde {f}) ^ {p} = \int_ {\mathrm{G/H}} ^ {*} \| \widetilde {f} \| ^ {p} d \nu = \int_ {\mathrm{G/H}} ^ {*} \gamma_ {\mathrm{G/H}} (g) (\| f \| ^ {p}) (y) \eta (g ^ {- 1}, y) d \nu (y)
\]

 以及 \((y\mapsto \eta (g^{-1},y))\cdot \nu = \pmb{\gamma}_{\mathrm{G / H}}(g)\nu\)，我们得到 \(\mathrm{N}_p(\widetilde{f}) = \mathrm{N}_p(f)\)。

 由这个引理可知，存在一个连续等距表示 \(\widetilde{\pi}\)，它从 G 映到 \(\mathrm{L}_{\overline{\pi}}^{p}(\mathrm{G})\) 中，使得对 \(f\in \mathcal{K}_{\pi}(\mathrm{G})\) 与 \(g\in \mathrm{G}\)，元素 \(\widetilde{\pi}(g)f\) 是上面定义的函数 \(\widetilde{f}\) 的类。若 \(p=2\)，则该表示是酉表示。

定义 2. --- 我们称如此定义的 G 在空间 \(L_{\overline{\pi}}^{2}(G)\) 中的酉表示为 G 的由 H 的表示 \(\pi\) 相对于 \(\kappa\) 所诱导的酉表示。把它记作 \(\operatorname{Ind}_{\mathrm{H}}^{\mathrm{G}}(\pi, \kappa)\)，或简作 \(\operatorname{Ind}_{\mathrm{H}}^{\mathrm{G}}(\pi)\)。

 注. --- 1) 设 \(\varrho\) 是 H 在复希尔伯特空间 F 中的一个酉表示，并设 \(u\colon\pi\to\varrho\) 是一个 H-态射。对每个函数 \(f\in\mathcal{K}_{\overline{\pi}}(\mathrm{G})\)，用 \(v(f)\) 表示 G 到 F 中的函数 \(g\mapsto u(f(g))\)；它属于 \(\mathcal{K}_{\overline{\varrho}}(\mathrm{G})\)，且满足 \(\mathrm{N}_{p}(v(f))\leqslant\|u\|\mathrm{N}_{p}(f)\)。由 \(\mathcal{K}_{\overline{\pi}}(\mathrm{G})\) 到 \(\mathcal{K}_{\overline{\varrho}}(\mathrm{G})\) 中的把 \(f\) 映为 \(v(f)\) 的线性映射因此延拓为 \(\mathrm{Ind}_{\mathrm{H}}^{\mathrm{G}}(\pi)\) 到 \(\mathrm{Ind}_{\mathrm{H}}^{\mathrm{G}}(\varrho)\) 中的一个连续 H-态射，记作 \(\mathrm{Ind}_{\mathrm{H}}^{\mathrm{G}}(u)\)。我们有 \(\mathrm{Ind}_{\mathrm{H}}^{\mathrm{G}}(1_{\pi})=1_{\mathrm{Ind}_{\mathrm{H}}^{\mathrm{G}}(\pi)}\)。设 \(\sigma\) 是 H 的一个酉表示，并设 \(v\colon\varrho\to\sigma\) 是一个 H-态射；我们有 \(\mathrm{Ind}_{\mathrm{H}}^{\mathrm{G}}(v\circ u)=\mathrm{Ind}_{\mathrm{H}}^{\mathrm{G}}(v)\circ\mathrm{Ind}_{\mathrm{H}}^{\mathrm{G}}(u)\)。

换言之，把 \(\pi\) 映为 \(\mathrm{Ind}_{\mathrm{H}}^{\mathrm{G}}(\pi)\)、把 \(u\) 映为 \(\mathrm{Ind}_{\mathrm{H}}^{\mathrm{G}}(u)\) 的构造是一个从 H 的酉表示的范畴到 G 的酉表示的范畴的函子（参见 CAT, I, § 2, 编纂中）。

\begin{enumerate}
\def\labelenumi{\arabic{enumi})}
\setcounter{enumi}{1}
\tightlist
\item
  设 \(\kappa^{\prime}\colon \mathrm{G}\to \mathbf{R}_{+}^{*}\)，使得
\end{enumerate}

\[
\frac {\kappa^ {\prime} (x h)}{\kappa^ {\prime} (x)} = \frac {\Delta_ {\mathrm{H}} (h)}{\Delta_ {\mathrm{G}} (h)} = \frac {\kappa (x h)}{\kappa (x)}
\]

 对所有 \((x,h)\in\mathrm{G}\times\mathrm{H}\) 成立。设 \(\nu^{\prime}\) 是 G/H 上的拟不变测度 \((\kappa^{\prime}\cdot\mu)/\beta\)。函数 \(\kappa^{\prime}\kappa^{-1}\) 通过取商定义了一个连续函数 \(\xi\colon G/H\to R_{+}^{*}\)，使得 \(\nu^{\prime}=\xi\cdot\nu\)。用 \(\alpha\) 表示 \(\mathcal{K}_{\pi}(G)\) 的由 \(f\mapsto(\kappa^{\prime}\kappa^{-1})^{1/p}f\) 定义的自同态；它满足

\[
\int_ {\mathrm{G} / \mathrm{H}} ^ {*} \| f \| ^ {p} d \nu^ {\prime} = \int_ {\mathrm{G} / \mathrm{H}} ^ {*} \| \alpha (f) \| ^ {p} d \nu
\]

并且它使我们可以把空间 \(\mathcal{L}_{\pi}^{p}(\mathrm{G},\nu^{\prime})\) 与 \(\mathcal{L}_{\pi}^{p}(\mathrm{G},\nu)\)、以及空间 \(\mathrm{L}_{\pi}^{p}(\mathrm{G},\nu^{\prime})\) 与 \(\mathrm{L}_{\pi}^{p}(\mathrm{G},\nu)\) 等同起来。此外，\(\alpha\) 诱导表示 \(\mathrm{Ind}_{\mathrm{H}}^{\mathrm{G}}(\pi,\kappa^{\prime})\) 与 \(\mathrm{Ind}_{\mathrm{H}}^{\mathrm{G}}(\pi,\kappa)\) 的一个等距同构。

\subsection{闭中心子群的情形}

在本目中，我们假设群 G 是单模的。

 考虑 G 的中心的一个闭子群 Z，并用 dz 表示 Z 上的一个哈尔测度。我们给商群 G/Z 赋以哈尔测度 \(\nu = \mu/dz\)。设 \(\chi\) 是 Z 的一个酉特征标。商群 G/Z 是单模的（INT, VII, 第 61 页, § 2, 第 7 目, 命题 11 的推论）。

引理 7. --- 我们有

\[
\mathrm{N} _ {1} (\check {f}) = \mathrm{N} _ {1} (f), \quad \langle f _ {1} | f _ {2} \rangle = \langle \check {f} _ {1} | \check {f} _ {2} \rangle .\tag{6}
\]

对所有 \(f \in \mathcal{F}_{\chi}(\mathrm{G})\) 以及一切 \(f_1\) 与 \(f_2\)（均属于 \(\mathcal{L}_{\chi}^{2}(\mathrm{G})\)）成立。

这些公式都是定义的直接推论。

 对每个 \(g \in \mathbf{G}\) 与每个 \(f \in \mathcal{F}_{\chi}(\mathbf{G})\)，函数 \(x \mapsto f(g^{-1}x)\) 与 \(x \mapsto f(xg)\) 都属于 \(\mathcal{F}_{\chi}(\mathbf{G})\)。分别把它们记作 \(\boldsymbol{\gamma}_{\mathrm{G},\chi}(g)f\) 与 \(\boldsymbol{\delta}_{\mathrm{G},\chi}(g)f\)。映射 \(\boldsymbol{\gamma}_{\mathrm{G},\chi}\) 与 \(\boldsymbol{\delta}_{\mathrm{G},\chi}\) 是 \(\mathbf{G}\) 在 \(\mathcal{F}_{\chi}(\mathbf{G})\) 中的线性表示。对每个 \(z \in \mathbf{Z}\)，我们有

\[
\boldsymbol {\gamma} _ {\mathrm{G}, \chi} (g z) = \overline {{\chi (z)}} \boldsymbol {\gamma} _ {\mathrm{G}, \chi} (g), \quad \boldsymbol {\delta} _ {\mathrm{G}, \chi} (g z) = \chi (z) \boldsymbol {\delta} _ {\mathrm{G}, \chi} (g).\tag{7}
\]

设 \(f\in \mathcal{F}_{\chi}(\mathrm{G})\) 且 \(g\in \mathrm{G}\)；我们有

\[
| \boldsymbol {\gamma} _ {\mathrm{G}, \chi} (g) f | = \boldsymbol {\gamma} _ {\mathrm{G} / \mathrm{Z}} (g \mathrm{Z}) | f |, \quad | \boldsymbol {\delta} _ {\mathrm{G}, \chi} (g) f | = \boldsymbol {\delta} _ {\mathrm{G} / \mathrm{Z}} (g \mathrm{Z}) | f |,\tag{8}
\]

其中出现在这些等式中的所有函数都等同于 G/Z 上的函数。

 子空间 \(\mathcal{K}_{\chi}(\mathrm{G})\)（\(\mathcal{F}_{\chi}(\mathrm{G})\) 的）在 \(\gamma_{\mathrm{G},\chi}\) 与 \(\delta_{\mathrm{G},\chi}\) 下稳定。设 \(p\) 是一个实数 \(\geqslant 1\)。公式 (8) 蕴含：表示 \(\gamma_{\mathrm{G},\chi}\) 与 \(\delta_{\mathrm{G},\chi}\) 在限制到 \(\mathcal{K}_{\chi}(\mathrm{G})\) 上后可延拓为 G 到 \(\mathcal{L}_{\chi}^{p}(\mathrm{G})\) 中的连续等距线性表示，分别记作 \(\gamma_{\mathrm{G},\chi}^{(p)}\) 与 \(\delta_{\mathrm{G},\chi}^{(p)}\)。通过取商，这些表示也在 \(\mathrm{L}_{\chi}^{p}(\mathrm{G})\) 中定义了 G 的等距表示，仍以同样的记号记之。

 表示 \(\gamma_{\mathrm{G},\chi}^{(2)}\) 与 \(\delta_{\mathrm{G},\chi}^{(2)}\)（\(\mathrm{L}_{\chi}^{2}(\mathrm{G})\) 中的）是酉表示，将分别简记为 \(\gamma_{\mathrm{G},\chi}\) 与 \(\delta_{\mathrm{G},\chi}\)——当不致与 \(\mathcal{F}_{\chi}(\mathrm{G})\) 中的表示混淆时。我们也用 \(\varrho_{\mathrm{G},\chi}\) 表示 \(\mathrm{G} \times \mathrm{G}\) 在 \(\mathcal{L}_{\chi}^{2}(\mathrm{G})\) 或 \(\mathrm{L}_{\chi}^{2}(\mathrm{G})\) 中的由下式定义的连续表示

\[
\pmb {\varrho} _ {\mathrm{G}, \chi} (g, h) = \pmb {\gamma} _ {\mathrm{G}, \chi} (g) \circ \pmb {\delta} _ {\mathrm{G}, \chi} (h) = \pmb {\delta} _ {\mathrm{G}, \chi} (h) \circ \pmb {\gamma} _ {\mathrm{G}, \chi} (g)
\]

对所有 \((g,h)\in \mathrm{G}\times \mathrm{G}\) 成立（参见 V, 第 377 页, 引理 1）。

\(\gamma_{\mathrm{G},\chi}\) 在 \(\mathrm{L}_{\chi}^{2}(\mathrm{G})\) 中的表示不是别的，正是诱导表示 \(\mathrm{Ind}_{\mathrm{Z}}^{\mathrm{G}}(\overline{\chi})\)。

当 \(Z = \{e\}\) 时，由于 G 单模，下面的引理可由 INT, VIII, 第 166 页, § 4, 第 5 目, 命题 12 得出。

引理 8. --- 设 \(f_1 \in \mathcal{K}_\chi(\mathrm{G})\) 且 \(f_2 \in \mathcal{L}_\chi^2(\mathrm{G})\)。

\begin{enumerate}
\def\labelenumi{\alph{enumi})}
\item
  函数 \(f\)——它定义在 \(\mathrm{G}\) 上，由 \(f(g) = \langle f_1 | \boldsymbol{\gamma}_{\mathrm{G},\chi}(g)f_2 \rangle\)（对所有 \(g \in \mathrm{G}\)）给出——属于 \(\mathcal{L}_{\chi}^{2}(\mathrm{G})\)，且满足 \(\mathrm{N}_2(f) \leqslant \mathrm{N}_1(f_1)\mathrm{N}_2(f_2)\)；
\item
  函数 \(f\)——它定义在 \(G\) 上，由 \(f(g) = \langle f_1 | \delta_{G,\chi}(g)f_2 \rangle\)（对所有 \(g \in G\)）给出——属于 \(\mathcal{L}_{\chi}^{2}(G)\)，且满足 \(N_2(f) \leqslant N_1(f_1)N_2(f_2)\)。
\end{enumerate}

我们来证明 \(a\) )，断言 \(b\) ) 的证明是类似的。函数 \(f\) 连续，因此 \(\mu\) -可测。对每个 \(z \in \mathbf{Z}\) 与每个 \(g \in \mathbf{G}\)，我们有

\[
f (g z) = \langle f _ {1} | \gamma_ {\mathrm{G}, \chi} (g z) f _ {2} \rangle = \overline {{\chi (z)}} \langle f _ {1} | \gamma_ {\mathrm{G}, \chi} (g) f _ {2} \rangle
\]

（公式 (7), 第 417 页），因此 \(f \in \mathcal{F}_{\overline{\chi}}(\mathrm{G})\)。V 第 413 页的命题 5 蕴含现在只需证明 \(\mathrm{N}_2(f) \leqslant \mathrm{N}_1(f_1)\mathrm{N}_2(f_2)\)。

先设 \(f_2\) 属于 \(\mathcal{K}_{\chi}(\mathrm{G})\)。对每个 \(g \in \mathrm{G}\)，按定义有

\[
f (g) = \int_ {\mathrm{G/Z}} \overline {{f _ {1}}} \gamma_ {\mathrm{G}, \chi} (g) f _ {2} d \nu ,
\]

其中函数 \(\overline{f_1} \gamma_{\mathrm{G},\chi}(g)f_2\) 等同于 G/Z 上的一个函数。

我们定义一个函数 \(f_3\)（在 \(G\) 上）：令 \(f_3(g) = 0\)（若 \(f_1(g) = 0\)），否则令 \(f_3(g) = f_1(g)|f_1(g)|^{-1/2}\)。函数 \(f_3\) 属于 \(\mathcal{F}_{\chi}(G)\) 且满足 \(f_1 = |f_1|^{1/2}f_3\)；它是 \(\mu\) -可测且模 \(Z\) 在一个紧集外为零的，因为 \(f_1\) 如此。由于 \(|f_1|^{1/2} \in \mathcal{K}_1(G)\)，由此得出

\[
\overline {{f _ {1}}} \gamma_ {\mathrm{G}, \chi} (g) f _ {2} = | f _ {1} | ^ {1 / 2} \overline {{f _ {3}}} \gamma_ {\mathrm{G}, \chi} (g) f _ {2},
\]

\[
| \overline {{f _ {3}}} \gamma_ {\mathrm{G}, \chi} (g) f _ {2} | = | f _ {3} | | \gamma_ {\mathrm{G}, \chi} (g) f _ {2} |
\]

对所有 \(g \in \mathbf{G}\) 成立。

 设 \(g \in G\)。由柯西--施瓦茨不等式，我们有

\[
\begin{array}{l} | f (g) | ^ {2} \leqslant \Big (\int_ {\mathrm{G/Z}} | f _ {1} | ^ {1 / 2} | f _ {3} | | \boldsymbol {\gamma} _ {\mathrm{G}, \chi} (g) f _ {2} | d \nu \Big) ^ {2} \\ \leqslant \Big (\int_ {\mathrm{G/Z}} | f _ {1} | d \nu \Big) \Big (\int_ {\mathrm{G/Z}} | f _ {3} | ^ {2} | \boldsymbol {\gamma} _ {\mathrm{G}, \chi} (g) f _ {2} | ^ {2} d \nu \Big). \end{array}
\]

由于有 \(|f_{3}|^{2}=|f_{1}|\)（在 \(\mathcal{K}(\mathrm{G}/\mathrm{Z})\) 中），我们通过在 G/Z 上积分得到

\[
\int_ {\mathrm{G} / \mathrm{Z}} | f (g) | ^ {2} d \nu (g) \leqslant \mathrm{N} _ {1} (f _ {1}) \int_ {\mathrm{G} / \mathrm{Z}} \left(\int_ {\mathrm{G} / \mathrm{Z}} | f _ {1} (x) | | f _ {2} (g ^ {- 1} x) | ^ {2} d \nu (x)\right) d \nu (g).
\]

由函数

\[
(g, x) \mapsto | f _ {1} (x) | | f _ {2} (g ^ {- 1} x) | ^ {2}
\]

通过取商所得到的 G/Z × G/Z 上的函数是 \((\nu \otimes \nu)\) 可测且具紧支撑的，从而是 \((\nu \otimes \nu)\) -缓增的（INT, V, 第 4 页, § 5, 第 2 目, 定义 2）。由 INT, V, 第 93 页, § 8, 第 3 目, 命题 7 得

\[
\begin{array}{l} \int_ {\mathrm{G} / \mathrm{Z}} \Bigl (\int_ {\mathrm{G} / \mathrm{Z}} | f _ {1} (x) |   | f _ {2} (g ^ {- 1} x) | ^ {2} d \nu (x) \Bigr) d \nu (g) \\ = \int_ {\mathrm{G} / \mathrm{Z}} | f _ {1} (x) | \Bigl (\int_ {\mathrm{G} / \mathrm{Z}} | f _ {2} (g ^ {- 1} x) | ^ {2} d \nu (g) \Bigr) d \nu (x) = \mathrm{N} _ {1} (f _ {1}) \mathrm{N} _ {2} (f _ {2}) ^ {2}. \end{array}
\]

于是我们有 \(\mathrm{N}_{2}(f)^{2}\leqslant\mathrm{N}_{1}(f_{1})^{2}\mathrm{N}_{2}(f_{2})^{2}\)，这就在此情形下建立了所要的性质。

 我们来考虑一般情形。用 \(u\) 表示 \(\mathcal{K}_{\chi}(\mathrm{G})\) 到 \(\mathcal{L}_{\overline{\chi}}^{2}(\mathrm{G})\) 中的把 \(f_{2}\) 映为 \(f\) 的线性映射。设 \(f_{2} \in \mathcal{L}_{\chi}^{2}(\mathrm{G})\)，并设 \((f_{2,n})_{n \in \mathbf{N}}\) 是 \(\mathcal{K}_{\chi}(\mathrm{G})\) 中的收敛于 \(f_{2}\)（在 \(\mathcal{L}_{\chi}^{2}(\mathrm{G})\) 中）的序列。令 \(f_{n} = u(f_{2,n})\)；序列 \((f_{n})_{n \in \mathbf{N}}\) 在 \(\mathcal{L}_{\overline{\chi}}^{2}(\mathrm{G})\) 中是柯西序列，因为前一种情形蕴含 \(\mathrm{N}_{2}(f_{n} - f_{m}) \leqslant \mathrm{N}_{1}(f_{1})\mathrm{N}_{2}(f_{2,n} - f_{2,m})\) 对一切 \(n\) 与 \(m\)（\(\mathbf{N}\) 中的）成立。设 \(f \in \mathcal{L}_{\overline{\chi}}^{2}(\mathrm{G})\) 使得 \((f_{n})\) 收敛于 \(f\)（V, 第 409 页, 命题 3, \(c\) )）。由于 \(\mathrm{N}_{2}(f_{n}) \leqslant \mathrm{N}_{1}(f_{1})\mathrm{N}_{2}(f_{2,n})\) 对所有 \(n \in \mathbf{N}\) 成立，由此得出 \(\mathrm{N}_{2}(f) \leqslant \mathrm{N}_{1}(f_{1})\mathrm{N}_{2}(f_{2})\)。

存在一个子序列 \((f_{n_{k}})_{k\in\mathbf{N}}\)，使得 \(f_{n_{k}}(g)\) 收敛于 \(f(g)\)，对所有位于 G 中模 Z 零测子集之外的 \(g\in G\) 成立（出处同上, b)）。但另一方面，对每个 \(g\in G\)，我们有

\[
f _ {n _ {k}} (g) = \left\langle f _ {1} \mid \gamma_ {\mathrm{G}, \chi} (g) f _ {2, n _ {k}} \right\rangle\rightarrow \left\langle f _ {1} \mid \gamma_ {\mathrm{G}, \chi} (g) f _ {2} \right\rangle .
\]

于是我们有 \(f(g) = \langle f_1 | \gamma_{\mathrm{G},\chi}(g)f_2 \rangle\)，对所有位于 G 中模 Z 零测子集之外的 \(g\) 成立。由于 \(f \in \mathcal{L}_{\overline{\chi}}^2(\mathrm{G})\) 且 \(\mathrm{N}_2(f) \leqslant \mathrm{N}_1(f_1)\mathrm{N}_2(f_2)\)，引理得证。

\subsection{平方可积表示}

在本目中，所考虑的希尔伯特空间都是复的。假设 G 是单模的，其中心记作 C。对 C 的每个闭子群 Z，用 \(\beta_{Z}\) 表示 Z 上的一个哈尔测度，并给 G/Z 赋以哈尔测度 \(\nu_{Z} = \mu/\beta_{Z}\)（INT, VII, 第 44 页, § 2, 第 2 目, 定义 1）。

设 \(\pi\) 是 G 在希尔伯特空间 E 中的一个不可约酉表示，\(\chi \in \widehat{\mathbf{C}}\) 是其中心特征标（V, 第 390 页, 定义 6）。对 E 中所有 \(x\) 与 \(y\)，用 \(f_{x,y}\) 表示矩阵系数 \(g \mapsto \langle x \mid \pi(g)y \rangle\)；它是 G 上取复值的连续函数。

设 E 中的 \(x\) 与 \(y\)。对 \(z \in \mathrm{C}\) 与 \(g \in \mathrm{G}\)，我们有

\[
f _ {x, y} (z g) = \langle x \mid \pi (z g) y \rangle = \langle x \mid \chi (z) \pi (g) y \rangle = \chi (z) f _ {x, y} (g),\tag{9}
\]

因此 \(f_{x,y}\) 属于空间 \(\mathcal{F}_{\chi}(\mathrm{G})\)（V, 第 408 页）。此外，对每对 \((g_1, h_1) \in \mathrm{G} \times \mathrm{G}\) 与每个 \(g \in \mathrm{G}\)，我们有

\[
f _ {\pi (g _ {1}) x, \pi (h _ {1}) y} (g) = \langle \pi (g _ {1}) x | \pi (g) \pi (h _ {1}) y \rangle = f _ {x, y} (g _ {1} ^ {- 1} g h _ {1}).\tag{10}
\]

关系式 (9) 为下面的定义提供了依据。

定义 3. --- 设 \(\pi\) 是 G 在希尔伯特空间 E 中的一个不可约酉表示。我们称 \(\pi\) 是模中心平方可积的，如果由函数 \(|f_{x,y}|\) 通过取商得到的 G/C 上的函数属于 \(\mathcal{L}^2 (\mathrm{G / C})\)，对所有 \((x,y)\in \mathrm{E}\times \mathrm{E}\) 成立。

这一条件不依赖于 G/C 上哈尔测度的选取。

若 \(\pi\) 的矩阵系数都属于 \(\mathcal{L}^2 (\mathrm{G})\)，则称 \(\pi\) 是平方可积的；G 的一个平方可积不可约酉表示的存在蕴含 G 的中心是紧的（V, 第 487 页, 习题 5）。

存在这样的群 G，它没有任何平方可积表示，即使模中心也没有（参见 V, 第 516 页, 习题 32）。

命题 6. --- 设 \(\pi\) 是 G 在希尔伯特空间 E 中的一个不可约酉表示。设 \(\chi\) 是 \(\pi\) 的中心特征标。则 \(\pi\) 模中心平方可积当且仅当存在 E 中非零元素 \(x_0\) 与 \(y_0\)，使得由 \(|f_{x_0,y_0}|\) 通过取商诱导的 G/C 上的函数属于 \(\mathcal{L}^2 (\mathrm{G / C})\)。

设存在 E 的非零元素 \(x_{0}\) 与 \(y_{0}\)，使得 \(|f_{x_{0},y_{0}}| \in \mathcal{L}^{2}(\mathrm{G/C})\)。只需证明这时 \(\pi\) 模中心平方可积。

 设 F 是由元素 \(x \in E\)——这里 \(|f_{x,y_{0}}|\) 属于 \(\mathcal{L}^{2}(G/C)\)——组成之集。这是 E 的一个向量子空间；它包含 \(x_{0}\)，从而非零。上面的关系式 (10) 蕴含 F 在 \(\pi\) 下稳定；由于表示 \(\pi\) 不可约，空间 F 在 E 中稠密。

 设 \(x \in \mathrm{F}\)。由于 \(f_{x,y_0}\) 属于 \(\mathcal{F}_{\chi}(\mathrm{G})\) 且 \(\mu\) -可测，又 \(|f_{x,y_0}| \in \mathcal{L}^2(\mathrm{G/C})\)，我们有 \(f_{x,y_0} \in \mathcal{L}_{\chi}^2(\mathrm{G})\)（V, 第 413 页的命题 5 应用于 Z = C）。用 \(u\) 表示 E 到 \(\mathrm{L}_{\chi}^{2}(\mathrm{G})\) 的部分算子，其定义域为 F，且把 \(x \in \mathrm{F}\) 映为 \(f_{x,y_0}\) 在 \(\mathrm{L}_{\chi}^{2}(\mathrm{G})\) 中的类。

 我们来证明部分算子 \(u\) 是闭的。设 \((x_{n}, u(x_{n}))_{n \in \mathbf{N}}\) 是 \(u\) 的图中的元素组成的序列，它在 \(\mathrm{E} \times \mathrm{L}_{\chi}^{2}(\mathrm{G})\) 中收敛。设 \(x\) 是 \((x_{n})\) 的极限，并设 \(f \in \mathcal{L}_{\chi}^{2}(\mathrm{G})\) 是一个函数，其类是序列 \((u(x_{n}))\) 的极限。

函数 \(u(x_{n})\) 是矩阵系数 \(f_{x_{n},y_{0}}\) 的类。对每个 \(g \in G\)，我们有

\[
f _ {x _ {n}, y _ {0}} (g) = \langle x _ {n} | \pi (g) y _ {0} \rangle \rightarrow \langle x | \pi (g) y _ {0} \rangle = f _ {x, y _ {0}} (g)
\]

当 \(n \to +\infty\)。因此我们有 \(f_{x,y_{0}} \in \mathcal{L}_{\chi}^{2}(\mathrm{G})\)，且模 C 几乎处处有 \(f = f_{x_{0},y}\)（V, 第 409 页, 命题 3 的推论）；这就是说 \(u(x)\) 是 f 在 \(\mathrm{L}_{\chi}^{2}(\mathrm{G})\) 中的类。于是我们证明了 u 是闭的。

u 的定义域 F 在 \(\pi\) 下稳定，且关系式 (10) 表明 u 满足 \(u \circ \pi(g) = \gamma_{\mathrm{G},\chi}(g) \circ u\)，对所有 \(g \in G\)。从而我们也有等式 \(u^{*} \circ \gamma_{\mathrm{G},\chi}(g) = \pi(g) \circ u^{*}\)，对所有 \(g \in G\)（V, 第 381 页, 引理 6），由此得 \((u^{*} \circ u) \circ \pi(g) = \pi(g) \circ (u^{*} \circ u)\)。而部分算子 \(u^{*} \circ u\) 是自伴的（IV, 第 241 页, 命题 12），特别地是闭的（IV, 第 236 页, 命题 7），于是 V, 第 387 页, 推论 1 蕴含 \(u^{*} \circ u\) 的定义域等于 E。更有，我们有 F = E，也就是说，函数 \(|f_{x,y_{0}}|\) 属于 \(\mathcal{L}^{2}(\mathrm{G}/\mathrm{C})\)，对所有 \(x \in E\) 成立。

设 \((x,y)\in\mathrm{E}\times\mathrm{E}\)。我们有 \(f_{x,y_{0}}\in\mathcal{L}_{\chi}^{2}(\mathrm{G})\)。把 \(\delta_{G,\chi}\) 代替 \(\gamma_{G,\chi}\)，用同样方法可证：\(y\in E\) 中使函数 \(|f_{x,y}|\) 属于 \(\mathcal{L}^{2}(\mathrm{G}/\mathrm{C})\) 的那种元素所成之集等于 E。命题得证。

在本节余下部分，我们取定 C 的一个使 C/Z 为紧的闭子群 Z。

引理 9. --- 设 \(\pi\) 是 G 在希尔伯特空间 E 中的一个模中心平方可积的不可约酉表示。设 \(\chi\) 是 \(\pi\) 的中心特征标到 Z 的限制。对 E 中所有 \(x\) 与 \(y\)，我们有 \(f_{x,y} \in \mathcal{L}_{\chi}^{2}(\mathrm{G})\)。

 函数 \(f_{x,y}\) 连续，因此 \(\mu\) -可测。由公式 (9)（第 420 页）我们有 \(f_{x,y} \in \mathcal{F}_{\chi}(\mathrm{G})\)。此外，由 INT, VII, 第 64 页, § 2, 第 8 目, 推论 1, c)，我们有 \(\mathrm{N}_{2}(f_{x,y}) < +\infty\)，因为 C/Z 紧且 \(|f_{x,y}| \in \mathcal{L}^{2}(\mathrm{G}/\mathrm{C})\)。于是断言由 V 第 413 页的命题 5 得出。

命题 7. --- 设 π 是 G 在希尔伯特空间 E 中的一个模中心平方可积的不可约酉表示。设 χ 是 π 的中心特征标到 Z 的限制。

存在一个实数 c \textgreater{} 0 与一个唯一的 \((\mathbf{G} \times \mathbf{G})\) -等距态射 w，它把酉表示 \(\overline{\pi} \boxtimes \pi\) 映到 \(\mathrm{L}_{\chi}^{2}(\mathrm{G})\) 中，使得对所有 \((x, y) \in \overline{\mathrm{E}} \times \mathrm{E}\)，元素 \(w(x \otimes y)\) 是 \(\mathrm{L}_{\chi}^{2}(\mathrm{G})\) 中函数 \(c^{1/2} f_{x,y}\) 的类。

 对每对 \((x,y)\in\mathrm{E}\times\mathrm{E}\)，我们有 \(f_{x,y}\in\mathcal{L}_{\chi}^{2}(\mathrm{G})\)（引理 9）。用 v 表示 \(\overline{E}\otimes E\) 到 \(\mathrm{L}_{\chi}^{2}(\mathrm{G})\) 中的唯一线性映射，使得 \(v(x\otimes y)\) 是 \(f_{x,y}\) 的类，对所有 \((x,y)\in\overline{\mathrm{E}}\times\mathrm{E}\) 成立。

我们在下面证明如下的引理。

引理 10. --- 存在一个实数 c \textgreater{} 0，使得线性映射 \(w = c^{1/2}v\) 是等距的。

假设该引理成立，注意公式 (10)（第 420 页）可以写成

\[
v (\overline {{\pi}} (g _ {1}) x \otimes \pi (h _ {1}) y) = \varrho_ {\mathrm{G}, \chi} (g _ {1}, h _ {1}) v (x \otimes y)\tag{11}
\]

 对所有 \((g_{1}, h_{1}) \in \mathbf{G} \times \mathbf{G}\) 与每对 \((x, y) \in \overline{\mathbf{E}} \times \mathbf{E}\) 成立。\(\overline{\mathbf{E}} \otimes \mathbf{E}\) 到 \(\mathrm{L}_{\chi}^{2}(\mathrm{G})\) 中的等距线性映射 w 可连续延拓（仍记作 w）到 \(\overline{E} \widehat{\otimes}_{2} E\) 上。由连续性与线性性，公式 (11) 蕴含 w 是一个 \((\mathrm{G} \times \mathrm{G})\) -态射——它从 \(\overline{\pi} \boxtimes \pi\) 映到 \(\mathrm{L}_{\chi}^{2}(\mathrm{G})\) 中——这就完成了命题的证明。

 我们来证明该引理。对所有 \(x \in E\)，用 \(u_{x}\) 表示线性映射 \(y \mapsto v(x \otimes y) = f_{x,y}\)（它从 E 映到 \(\mathrm{L}_{\chi}^{2}(\mathrm{G})\) 中）。我们有 \(u_{x} \in \mathrm{Hom}_{\mathrm{G}}(\pi, \gamma_{\mathrm{G},\chi})\)。由公式 (10)，根据 V 第 388 页的推论 5，存在一个实数 \(\lambda_{x} \geqslant 0\)，使得 \(\lambda_{x}u_{x}\) 是等距的。

设 E 中的 \(x\) 与 \(y\)。我们有

\[
\| f _ {x, y} \| ^ {2} = \int_ {\mathrm{G/Z}} \overline {{f}} _ {x, y} f _ {x, y} d \nu_ {\mathrm{Z}} = \int_ {\mathrm{G/Z}} \check {\overline {{f}}} _ {x, y} \check {f} _ {x, y} d \nu_ {\mathrm{Z}}
\]

因为 G/Z 是单模的（V 第 417 页, 引理 7）。注意到 \(\check{f}_{x,y} = \overline{f}_{y,x}\)，我们得到

\[
\lambda_ {x} \| y \| ^ {2} = \| f _ {x, y} \| ^ {2} = \int_ {\mathrm{G/Z}} f _ {y, x} \overline {{f}} _ {y, x} d \nu_ {\mathrm{Z}} = \| f _ {y, x} \| ^ {2} = \lambda_ {y} \| x \| ^ {2}.
\]

这表明正实数 \(\lambda_{x}/\|x\|^{2}\) 与 E 中非零元素 x 的选取无关。它是严格正的，因为对每个非零的 x，函数 \(f_{x,x}\) 连续且在 e 处取值 \(\|x\|^{2}>0\)，由此 \(\|u_{x}(x)\|=\|f_{x,x}\|>0\)。用 \(c^{-1}\) 记这个实数。

对每对 \((x,y)\in\mathrm{E}\times\mathrm{E}\)，由此得出

\[
\| v (x \otimes y) \| ^ {2} = \| f _ {x, y} \| ^ {2} = \lambda_ {x} \| y \| ^ {2} = c ^ {- 1} \| x \| ^ {2} \| y \| ^ {2} = c ^ {- 1} \| x \otimes y \| ^ {2}.
\]

利用 EVT, V, 第 29 页, 推论 1，我们推出线性映射 \(w = c^{1/2}v\)（它从 \(\overline{\mathrm{E}}\widehat{\otimes}_2\mathrm{E}\) 映到 \(\mathrm{L}_{\chi}^{2}(\mathrm{G})\) 中）是等距的，这正是所要证的。

 推论. --- 设 \(\pi\) 是 G 的一个不可约酉表示，\(\chi\) 是其中心特征标到 Z 的限制。表示 \(\pi\) 模中心平方可积当且仅当它同构于表示 \(\gamma_{\mathrm{G},\overline{\chi}}\)（从 G 到 \(\mathrm{L}_{\overline{\chi}}^{2}(\mathrm{G})\) 中）的一个子表示（相应地，同构于表示 \(\delta_{\mathrm{G},\chi}\)（从 G 到 \(\mathrm{L}_{\chi}^{2}(\mathrm{G})\) 中）的一个子表示）。

我们来证明关于 \(\gamma_{\mathrm{G},\overline{\chi}}\) 的断言，第二个断言的证明是类似的。

 先设表示 \(\gamma_{\mathrm{G},\overline{\chi}}\) 存在一个同构于 \(\pi\) 的子表示 E。由于空间 E 非零，存在函数 \(f_1 \in \mathcal{K}_{\overline{\chi}}(\mathrm{G})\)，其类 \(\widetilde{f}_1\) 不与 E 正交。我们有 \(f_1 \in \mathcal{L}_{\overline{\chi}}^1 (\mathrm{G})\)。用 \(\widetilde{f}_{1,\mathrm{E}}\) 表示 \(\widetilde{f}_1\) 在 E 上的正交投影；它是 E 的一个非零元素。再设非零的 \(\widetilde{f}_2 \in \mathrm{E}\)。映射 \(h: g \mapsto \langle \widetilde{f}_{1,\mathrm{E}} | \gamma_{\mathrm{G},\overline{\chi}}(g)\widetilde{f}_2 \rangle\) 是 \(\pi\) 的一个矩阵系数。对每个 \(g \in \mathrm{G}\)，我们有 \(h(g) = \langle \widetilde{f}_1 | \gamma_{\mathrm{G},\overline{\chi}}(g)\widetilde{f}_2 \rangle\)，因为 \(\widetilde{f}_1 - \widetilde{f}_{1,\mathrm{E}}\) 与 E 正交。由 V 第 418 页的引理 8，函数 \(h\) 属于 \(\mathcal{L}_{\chi}^{2}(\mathrm{G})\)，于是命题 6 蕴含 \(\pi\) 模中心平方可积。

反之，设 \(\pi\) 模中心平方可积。设 \(x_0 \in \mathrm{E}\) 是一个非零向量。由命题 7 与公式 (10)，映射 \(y \mapsto \check{f}_{x_0,y}\) 是 \(\pi\) 到 \(\gamma_{\mathrm{G},\overline{\chi}}\) 中的一个单射 G-态射。

定义 4. --- 设 Z 是 C 的一个使 C/Z 为紧的闭子群。设 π 是 G 的一个模中心平方可积的不可约酉表示。满足命题 7 之性质的唯一实数 c \textgreater{} 0 称为 π 相对于 Z 的形式次数。把它记作 \(c_{\mathrm{Z}}(\pi)\)。

形式次数依赖于 Z 上哈尔测度的选取。若把 Z 上的哈尔测度 \(\beta_{Z}\) 乘以一个实数 t \textgreater{} 0，则 G/Z 上的测度 \(\nu_{Z} = \mu/\beta_{Z}\) 被乘以 \(t^{-1}\)，而 \(\pi\) 的形式次数被乘以 t。

形式次数由下述性质所刻画：

命题 8（正交关系）. --- 设 π 是 G 在希尔伯特空间 E 中的一个模中心平方可积的不可约酉表示。则有

\[
c _ {\mathrm{Z}} (\pi) \int_ {\mathrm{G} / \mathrm{Z}} \overline {{\langle x | \pi (g) x ^ {\prime} \rangle}} \langle y | \pi (g) y ^ {\prime} \rangle d \nu_ {\mathrm{Z}} (g) = \overline {{\langle x | y \rangle}} \langle x ^ {\prime} | y ^ {\prime} \rangle
\]

对所有 \((x,y,x',y')\in \mathrm{E}^4\) 成立

用 \(w\) 表示命题 7 中的态射。我们有

\[
\int_ {\mathrm{G} / \mathrm{Z}} \overline {{\langle x | \pi (g) x ^ {\prime} \rangle}} \langle y | \pi (g) y ^ {\prime} \rangle d \nu_ {\mathrm{Z}} (g) = \langle f _ {x, x ^ {\prime}} | f _ {y, y ^ {\prime}} \rangle
\]

而根据出处同上，由此得出

\[
\begin{array}{c} \langle f _ {x, x ^ {\prime}} | f _ {y, y ^ {\prime}} \rangle = \frac {1}{c _ {\mathrm{Z}} (\pi)} \langle w (x \otimes x ^ {\prime}) | w (y \otimes y ^ {\prime}) \rangle \\ = \frac {1}{c _ {\mathrm{Z}} (\pi)} \langle x \otimes x ^ {\prime} | y \otimes y ^ {\prime} \rangle = \frac {1}{c _ {\mathrm{Z}} (\pi)} \langle x | y \rangle_ {\overline {{\mathrm{E}}}} \langle x ^ {\prime} | y ^ {\prime} \rangle_ {\mathrm{E}}, \end{array}
\]

由此即得结论。

除上述命题外，对于互不同构的平方可积不可约表示，我们还有下述关系。

命题 9. --- 设 \(\pi_1\) 与 \(\pi_2\) 是 G 在希尔伯特空间 E \(_1\) 与 E \(_2\) 中的两个互不同构的不可约酉表示。假设 \(\pi_1\) 与 \(\pi_2\) 都模中心平方可积，且它们的中心特征标到 Z 的限制相同。则有

\[
\int_ {\mathrm{G} / \mathrm{Z}} \overline {{\langle x \mid \pi_ {1} (g) x ^ {\prime} \rangle}} \langle y \mid \pi_ {2} (g) y ^ {\prime} \rangle d \nu_ {\mathrm{Z}} (g) = 0
\]

对所有 \((x,x^{\prime},y,y^{\prime})\in \mathrm{E}_{1}^{2}\times \mathrm{E}_{2}^{2}\) 成立

 对 \(i=1,2\)，用 \(w_{i}\) 表示关于表示 \(\pi_{i}\) 的命题 7 中的态射。由 V 第 384 页的引理 8 \(b\) ) 与 V 第 394 页的命题 8 的断言 \(b\) )，\(w_{i}\) 的像含于 \(\pi_{i}\) -同型分量（\(\delta_{\mathrm{G},\chi}\) 的）中。由出处同上的断言 \(a\) )，\(w_{1}\) 的像因此与 \(w_{2}\) 的像正交。于是得出 \(\langle w_{1}(x\otimes x^{\prime})|w_{2}(y\otimes y^{\prime})\rangle=0\)，对所有 \((x,x^{\prime},y,y^{\prime})\in\mathrm{E}_{1}^{2}\times\mathrm{E}_{2}^{2}\) 成立，这就是所要证的公式。

注. --- 命题 8 与命题 9 的正交关系推广了 A, VIII, 第 399 页中的正交关系（另见 V 第 457 页 § 2 中 G 为紧的情形）。

\subsection{交换群正则表示的子表示}

在本节中，G 是一个局部紧交换群，\(\mu\) 是 G 上的一个哈尔测度。用 \(\hat{G}\) 表示 G 的对偶群（II, 第 201 页, 定义 2），用 \(\hat{\mu}\) 表示 \(\mu\) 在 \(\hat{G}\) 上的对偶哈尔测度（II, 第 214 页, 定义 4）。可测性概念总是相对于 \(\mu\) 与 \(\hat{\mu}\) 而言。

我们打算确定左正则表示 \(\gamma_{G}\)（G 在 \(L^{2}(G,\mu)\) 中的）的所有子表示。由于 G 是交换的，我们还有 \(\delta_{\mathrm{G}}(g)=\gamma_{\mathrm{G}}(g^{-1})\)，故这些子表示也是右正则表示的子表示。

对 \(\widehat{\mathbf{G}}\) 的每个可测子集 M，用 \(\mathrm{E_M}\) 表示满足下列条件的 \(f \in \mathrm{L}^2 (\mathrm{G},\mu)\) 组成的集合：傅里叶变换 \(\mathcal{F}_{\mathrm{G}}(f)\) 在 M 上几乎处处为零（参见 II, 第 210 页, 第 3 目）。它是连续线性映射 \(f \mapsto \varphi_{\mathrm{M}}\mathcal{F}_{\mathrm{G}}(f)\)（从 \(\mathrm{L}^2 (\mathrm{G},\mu)\) 到 \(\mathrm{L}^2 (\widehat{\mathrm{G}},\widehat{\mu})\) 中）的核，其中 \(\varphi_{\mathrm{M}}\) 表示 M 的特征函数（参见 II, 第 215 页, 定理 1），因而是 \(\mathrm{L}^2 (\mathrm{G},\mu)\) 的一个闭子空间。

若 \(\widehat{G}\) 的可测子集 M 与 N 使得 \((\mathrm{M} \cup \mathrm{N}) - (\mathrm{M} \cap \mathrm{N})\) 是局部零测的，则称它们相差一个局部零测集意义下相等。等价地，可测子集 M 与 N 相差一个局部零测集意义下相等，当且仅当 M 与 N 的特征函数在 \(\mathrm{L}^{\infty}(\widehat{\mathrm{G}}, \widehat{\mu})\) 中相等。

命题 10. --- a) 设 M 是 \(\widehat{G}\) 的一个可测子集。空间 \(\mathrm{E}_{\mathrm{M}}\) 是表示 \(\gamma_{\mathrm{G}}\) 的一个子表示；

\begin{enumerate}
\def\labelenumi{\alph{enumi})}
\setcounter{enumi}{1}
\item
  设 M 与 N 是 \(\widehat{\mathbf{G}}\) 的可测子集。我们有 \(\mathrm{E_M} = \mathrm{E_N}\) 当且仅当 M 与 N 相差一个局部零测集意义下相等；
\item
  \(\gamma_{\mathrm{G}}\) 的任何子表示都形如 \(\mathrm{E}_{\mathrm{M}}\)，其中 M 是 \(\widehat{\mathrm{G}}\) 的一个可测子集。
\end{enumerate}

用 \(\eta\) 表示 G 到 \(\widehat{\mathbf{G}}\) 的典范映射（参见 II, 第 216 页, 注 1）。由于 \(\mathrm{E_M}\) 是闭的，断言 \(a\) ) 由以下公式得出

\[
\boldsymbol {\gamma} _ {\mathrm{G}} (x) (f) = \varepsilon_ {x} * f, \quad \mathcal {F} _ {\mathrm{G}} (\varepsilon_ {x} * f) = \eta (x) \mathcal {F} _ {\mathrm{G}} (f)
\]

其中 \(x\in \mathrm{G}\)，\(f\in \mathrm{L}^2 (\mathrm{G},\mu)\)。

设 M 与 N 是相差一个局部零测集意义下相等的可测子集。由于此时 \(\varphi_{M}\) 等于 \(\varphi_{N}\)（在 \(L^{\infty}(\widehat{G},\widehat{\mu})\) 中），这一条件蕴含 \(E_{M}=E_{N}\)。

反之，设 M 与 N 不是相差一个局部零测集意义下相等的。那么在交换 M 与 N 的角色之后，存在紧子集 K 使得集合 \(L = K \cap (M - (M \cap N))\) 不是零测的。设 \(\varphi_{\mathrm{L}} \in \mathrm{L}^{2}(\widehat{\mathrm{G}}, \widehat{\mu})\) 是 L 的特征函数所在的类，并令 \(f = \overline{\mathcal{F}}_{\widehat{\mathrm{G}}}(\varphi_{\mathrm{L}}) \in \mathrm{L}^{2}(\mathrm{G}, \mu)\)；我们有 \(\mathcal{F}_{\mathrm{G}}(f) = \varphi_{\mathrm{L}}\)（II, 第 220 页, 定理 2 的推论）。于是 f 属于 \(E_{N}\)，因为 \(L \cap N\) 是空集；但不属于 \(E_{M}\)，因为 \(\varphi_{\mathrm{M}}\mathcal{F}(f) = \varphi_{\mathrm{M}}\varphi_{\mathrm{L}} = \varphi_{\mathrm{L}}\)。因此 \(E_{M} \neq E_{N}\)，这就证明了 b)。

现在设 E 是 \(\gamma_{\mathrm{G}}\) 的一个子表示。设 \(p_{\mathrm{E}}\) 是 \(\mathrm{L}^2 (\mathrm{G},\mu)\) 中以 E 为像的正交投影算子，并令 \(q_{\mathrm{E}} = \mathcal{F}_{\mathrm{G}}\circ p_{\mathrm{E}}\circ \mathcal{F}_{\mathrm{G}}^{-1}\)。投影算子 \(p_{\mathrm{E}}\) 属于 \(\mathrm{Hom}_{\mathrm{G}}(\gamma_{\mathrm{G}},\gamma_{\mathrm{G}})\)（V, 第 383 页, 命题 4），故它与 \(\gamma_{\mathrm{G}}(f)\) 交换，对所有 \(f\in \mathrm{L}^{1}(\mathrm{G},\mu)\)（参见 V, 第 401 页）。这意味着它与自同态 \(\varphi \mapsto f*\varphi\) 交换，对所有 \(f\in \mathrm{L}^{1}(\mathrm{G},\mu)\)。于是由 V 第 407 页的引理 4，自同态 \(q_{\mathrm{E}}\)（\(\mathrm{L}^2 (\widehat{\mathrm{G}},\widehat{\mu})\) 的）与乘以 \(g\) 的乘法自同态交换，对每个满足下列条件的函数 \(g\in \mathcal{C}_0(\widehat{\mathrm{G}})\)：它属于 \(\mathrm{L}^1 (\mathrm{G},\mu)\) 的傅里叶变换在 \(\mathcal{C}_0(\widehat{\mathrm{G}})\) 中的像（II, 第 223 页, 命题 14）。由于傅里叶变换的像在 \(\mathcal{C}_0(\widehat{\mathrm{G}})\) 中稠密（II, 第 209 页, 命题 5 的推论），态射 \(g\mapsto m_g\) 的连续性蕴含 \(q_{\mathrm{E}}\) 与 \(m_g\) 交换，对每个函数 \(g\in \mathcal{C}_0(\widehat{\mathrm{G}})\)。

于是由 IV 第 188 页的命题 7 的推论，存在函数 \(\varphi \in \mathcal{L}^{\infty}(\widehat{\mathrm{G}},\widehat{\mu})\) 使得 \(q_{\mathrm{E}} = m_{\varphi}\)。由于 \(q_{\mathrm{E}} = q_{\mathrm{E}}^{2}\)，我们有 \(m_{\varphi} = m_{\varphi}^{2} = m_{\varphi^{2}}\)，从而 \(\varphi = \varphi^{2}\)（在 \(\mathrm{L}^{\infty}(\widehat{\mathrm{G}},\widehat{\mu})\) 中）（IV, 第 186 页, 命题 5）；这意味着 \(\varphi\) 在 \(\mathrm{L}^{\infty}(\widehat{\mathrm{G}},\widehat{\mu})\) 中等于 \(\widehat{\mathrm{G}}\) 的某个可测子集 N 的特征函数所在的类。令 \(\mathrm{M} = \widehat{\mathrm{G}} - \mathrm{N}\)。设 \(f \in \mathrm{L}^{2}(\mathrm{G},\mu)\)。我们有：\(f \in \mathrm{E}\) 当且仅当 \(p_{\mathrm{E}}(f) = f\)，当且仅当 \(\varphi \mathcal{F}_{\mathrm{G}}(f) = \mathcal{F}_{\mathrm{G}}(f)\)，而这等价于 \(f \in \mathrm{E}_{\mathrm{M}}\)。

注记. --- 设 \(\chi\) 是 G 的一个特征标。若 G 不是紧的，则 \(\chi\) -同型分量（G 在 L \(^{2}\) (G) 中的正则表示的）

是平凡的。若 G 是紧的，则 G 的正则表示的 \(\chi\) -同型分量的维数为 1，且函数 \(\chi \in \mathrm{L}^{2}(\mathrm{G})\) 是它的一个基。

\subsection{群 R 的酉表示}

在本节中，E 表示一个复希尔伯特空间。群 R 赋以勒贝格测度。

引理 11. --- 设 u 是 E 上的一个自伴部分算子。对 \(t \in \mathbf{R}\)，令 \(\varrho(t) = e^{itu} \in \mathcal{L}(\mathrm{E})\)。则映射 \(\varrho\) 是 \(\mathbf{R}\) 的一个酉表示。

算子 \(\varrho(t)\) 由普遍可测函数演算定义（IV, 第 272 页, 定义 5）；它是 E 的一个自同态，因为函数 \(x \mapsto e^{itx}\) 在 \(\mathrm{Sp}(u)\) 上有界（IV, 第 275 页, 命题 5, \(a\) )）。

由 IV 第 275 页的命题 5，我们有 \(\varrho(0) = 1_{\mathrm{E}}\)，\(\varrho(t)^* = \varrho(-t)\) 对所有 \(t \in \mathbf{R}\)，且 \(\varrho(t_1 + t_2) = \varrho(t_1)\varrho(t_2)\) 对所有 \(t_1\) 与 \(t_2\)（都属于 \(\mathbf{R}\)）。特别地，自同态 \(\varrho(t)\) 对每个 \(t \in \mathbf{R}\) 都是酉的。对每个 \(x \in \mathrm{E}\)，从 \(\mathbf{R}\) 到 E 的由 \(x \mapsto \varrho(t)x\) 定义的映射在 \(t = 0\) 处连续（IV, 第 276 页, 命题 6），故 \(\varrho\) 是 \(\mathbf{R}\) 在 E 中的一个酉表示（V, 第 380 页, 引理 4）。

引理 12. --- 设 \(\varrho\) 是 \(\mathbf{R}\) 在 E 中的一个酉表示。设 D 是 E 中由满足下列条件的元素 \(x\) 组成的集合：映射 \(\psi_x: t \mapsto \varrho(t)x\) 在 0 处可微。D 到 E 中的映射 \(u\) 由 \(x \mapsto i^{-1} \psi_x'(0)\) 给出，它定义了 E 上一个对称的部分算子。

设 \(f \in \mathcal{D}(\mathbf{R})\) 且 \(x \in \mathrm{E}\)。令 \(y = \varrho(f)x\)。对每个 \(h \in \mathbf{R}\)，我们有

\[
\begin{array}{r l} & {\psi_ {y} (h) = \varrho (h) \varrho (f) x = \varrho (f) \varrho (h) x} \\ & {\qquad = \int_ {\mathbf {R}} f (t) \varrho (t + h) x d t = \int_ {\mathbf {R}} f (t - h) \psi_ {x} (t) d t.} \end{array}
\]

从 \(R^{2}\) 到 C 的由 \((t,h)\mapsto f(t-h)\) 定义的函数是无穷次可微的；它对 h 的导数是由 \((t,h)\mapsto -f'(t-h)\) 定义的函数，后者被 \(\|f'\|_{\infty}\) 所控制。当 h=0 时，这个导数在 t 属于一个紧集之外时为零。由于对所有 t 有 \(\|\psi_{x}(t)\|=\|x\|\)，由 IV 第 197 页的命题 2 推出，函数 \(\psi_{y}\) 在 0 处可微且其导数由下式给出

\[
\psi_ {y} ^ {\prime} (0) = - \int_ {\mathbf {R}} f ^ {\prime} (t) \psi_ {x} (t) d t = - \varrho (f ^ {\prime}) x.
\]

因此 \(y \in D\)。由于空间 \(\mathcal{D}(\mathbf{R})\) 在 \(\mathrm{L}^{1}(\mathbf{R})\) 中稠密，由 INT, VIII, 第 139 页, § 2, 第 7 目, 命题 10 以及（IV, 第 202 页, 命题 4）推出，空间 D 在 E 中稠密。

设 \(x_{1}\) 与 \(x_{2}\) 属于 D。我们计算

\[
\begin{array}{c} \langle u (x _ {1}) \mid x _ {2} \rangle = \lim _ {h \to 0} \frac {i}{h} \langle (\varrho (h) - 1 _ {\mathrm{E}}) x _ {1} \mid x _ {2} \rangle \\ = \lim _ {h \to 0} \frac {i}{h} \langle x _ {1} \mid (\varrho (- h) - 1 _ {\mathrm{E}}) x _ {2} \rangle = \langle x _ {1} \mid u (x _ {2}) \rangle . \end{array}
\]

因此，部分算子 \(u\) 是对称的。

定义 5. --- 设 \(\varrho\) 是 \(\mathbf{R}\) 在 E 中的一个酉表示。引理 12 中定义的对称部分算子称为 \(\varrho\) 的无穷小生成元。

定理 1（斯通）. --- 映射 σ 把 R 在 E 中的每个酉表示 ρ 对应到其无穷小生成元 u，它定义了从 R 在 E 中的酉表示集合到 E 上自伴部分算子集合上的一个双射。逆双射 τ 把一个自伴部分算子对应到酉表示 \(t \mapsto e^{itu}\)。

我们先证明几个引理。

引理 13. --- 设 \(\varrho\) 是 \(\mathbf{R}\) 在 E 中的一个酉表示，\(u\) 是其无穷小生成元。

\begin{enumerate}
\def\labelenumi{\alph{enumi})}
\item
  u 的定义域是满足下列条件的 \(x \in E\) 组成的集合：映射 \(\psi_{x} \colon t \mapsto \varrho(t)x\) 在 R 上可微；
\item
  对每个 \(t \in \mathbf{R}\) 与每个 \(x \in \mathrm{dom}(u)\)，我们有 \(\varrho(t)x \in \mathrm{dom}(u)\) 且 \(\psi_x'(t) = i u(\varrho(t)x)\)；
\item
  部分算子 u 是本质自伴的。
\end{enumerate}

对每个 \(x \in \mathrm{E}\)，用 \(\psi_x\) 表示从 \(\mathbf{R}\) 到 \(\mathrm{E}\) 中的由 \(\psi_x(t) = \varrho(t)x\) 定义的映射。对每个 \(t \in \mathbf{R}\) 与每个 \(h \in \mathbf{R}\)，我们有

\[
\psi_ {x} (t + h) - \psi_ {x} (t) = \varrho (t) \big (\psi_ {x} (h) - \psi_ {x} (0) \big),
\]

这证明了 \(\operatorname{dom}(u)\) 是满足下列条件的元素 \(x \in E\) 组成的空间：\(\psi_{x}\) 在 R 上可微；并且确立了 \(\psi_{x}^{\prime}(t) = \varrho(t)\psi_{x}^{\prime}(0) = i\varrho(t)(u(x))\) 对所有 \(t \in R\)。

设 \(x \in E\) 且 \(t \in R\)。我们有 \(\psi_{\varrho(t)x}(s) = \psi_x(s + t)\)，对所有 \(s \in R\)。于是我们有 \(\varrho(t)x \in \text{dom}(u)\)，只要 \(x \in \text{dom}(u)\)，而且由 a) 有 \(u(\varrho(t)x) = i^{-1}\psi_x'(t) = \varrho(t)u(x)\)。由此得到断言 b)。

我们来证明 c)。设 \(x \in \text{Ker}(u^{*} - i1_{\text{E}})\)。我们来证明 x = 0。设 \(y \in \text{dom}(u)\)，并设 f 是 R 上由 \(f(t) = \langle \psi_{y}(t) | x \rangle\)（\(t \in R\)）定义的函数。函数 f 是有界的，因为 \(\| \psi_{y}(t) \| = \| y \|\) 对所有 \(t \in R\)；它在 R 上可微，且对每个 \(t \in R\)，断言 b) 蕴含

\[
\begin{array}{r l} f ^ {\prime} (t) = \langle \psi_ {y} ^ {\prime} (t) | x \rangle & = \langle i u (\varrho (t) y) | x \rangle \\ & = - i \langle \varrho (t) y | u ^ {*} (x) \rangle = - i \langle \psi_ {y} (t) | i x \rangle = f (t) \end{array}
\]

因为 \(u^{*}(x)=ix\)。因此我们有 \(f(t)=f(0)e^{t}\)，对所有 \(t\in R\)（FVR, IV, 第 27 页）。由于 f 有界，函数 f 为零，特别地我们有 \(\langle y|x\rangle=f(0)=0\)。由于空间 dom(u) 在 E 中稠密，我们得出 x=0。类似地，我们证明 \(\operatorname{Ker}(u^{*}+i1_{\mathrm{E}})\) 化为 0。由 IV 第 261 页的推论 3，部分算子 u 因此是本质自伴的。

引理 14. --- 设 u 是 E 上的一个自伴部分算子，\(\varrho(t) = e^{itu}\) 是由 u 定义的 R 的酉表示。\(\varrho\) 的无穷小生成元等于 u。

设 \(x \in \text{dom}(u)\)。令 \(\psi_{x}(t) = \varrho(t)x\)，对所有 \(t \in R\)。对每个非零实数 h，我们有

\[
\frac {1}{h} \left(\psi_ {x} (t + h) - \psi_ {x} (t)\right) = \left(\frac {1}{h} \left(e ^ {i h u} - 1 _ {\mathrm{E}}\right)\right) e ^ {i t u} x = \left(\frac {1}{h} \left(e ^ {i h u} - 1 _ {\mathrm{E}}\right)\right) \varrho (t) x.
\]

当 h 趋于 0 时，我们有

\[
\frac {1}{h} (e ^ {i h t} - 1) \rightarrow i t
\]

对所有 \(t \in R\) 成立。此外，

\[
\left| \frac {1}{h} (e ^ {i h t} - 1) \right| = | t | \left| \frac {1}{h} \int_ {0} ^ {h} e ^ {i t s} d s \right| \leqslant | t |.
\]

于是由 IV 第 276 页的命题 6 推出，函数 \(\psi_x\) 在 \(\mathbf{R}\) 上可微且满足 \(\psi_x'(t) = i u(\varrho(t)x)\)，对所有 \(t \in \mathbf{R}\)。因此，\(u\) 的定义域含于满足下列条件的 \(x \in E\) 组成的集合：\(\psi_x\) 在 0 处可微且此时有 \(\psi_x'(0) = i u(x)\)。按定义，这意味着 \(\varrho\) 的无穷小生成元是 \(u\) 的一个扩张。由于这两个算子都是对称的而 \(u\) 是自伴的（IV, 第 238 页, 注 5），它们因此相等。

引理 15. --- 设 \(\varrho\) 是 \(\mathbf{R}\) 在 E 中的一个酉表示。则无穷小生成元 \(u\)（\(\varrho\) 的）是自伴的，且 \(\varrho(t) = e^{itu}\) 对所有 \(t \in \mathbf{R}\)。

由引理 13(c)，部分算子 u 是本质自伴的。因此它的闭包 \(\overline{u}\) 是一个自伴算子。用 \(\pi\) 表示由 \(\pi(t) = e^{it\overline{u}}\) 定义的 R 的酉表示（引理 11）。

对每个 \(x \in \mathrm{E}\)，用 \(\psi_x\)（相应地 \(\widetilde{\psi}_x\)）表示从 \(\mathbf{R}\) 到 \(\mathrm{E}\) 中的由 \(\psi_x(t) = \varrho(t)x\)（相应地由 \(\widetilde{\psi}_x(t) = \pi(t)x\)）定义的映射。

由引理 13 的 a) 与 b) 两部分，空间 \(\mathrm{dom}(u)\) 是 E 中由满足下列条件的元素 \(x \in \mathrm{E}\) 组成的子空间：映射 \(\psi_x\) 在 \(\mathbf{R}\) 上可微，且对所有 \(x \in \mathrm{dom}(u)\) 与每个 \(t \in \mathbf{R}\)，我们有 \(\psi_x'(t) = i u (\psi_x(t))\)。同样，空间 \(\mathrm{dom}(\overline{u})\) 是 E 中由满足下列条件的元素 \(x \in \mathrm{E}\) 组成的子空间：映射 \(\widetilde{\psi}_x\) 在 \(\mathbf{R}\) 上可微，且对所有 \(x \in \mathrm{dom}(\overline{u})\) 与每个 \(t \in \mathbf{R}\)，我们有 \(\widetilde{\psi}_x'(t) = i \overline{u} (\widetilde{\psi}_x(t))\)。

设 \(x \in \text{dom}(u) \subset \text{dom}(\overline{u})\)。令 \(f = \psi_x - \widetilde{\psi}_x\)。这是从 \(\mathbf{R}\) 到 E 中的一个可微函数。对每个 \(t \in \mathbf{R}\)，由此得出

\[
f ^ {\prime} (t) = i u (\psi_ {x} (t)) - i \overline {{u}} (\widetilde {\psi} _ {x} (t)) = i \overline {{u}} (f (t))
\]

因为 \(\psi_x(t) \in \mathrm{dom}(u)\) 且 \(u \subset \overline{u}\)。令 \(g = \|f\|^2\)；它是从 \(\mathbf{R}\) 到 \(\mathbf{R}\) 中的一个可微映射且满足 \(g(0) = 0\)。对 \(t \in \mathbf{R}\)，由 FVR, I, 第 28 页, 命题 2 得到

\[
\begin{array}{r l} g ^ {\prime} (t) & = \langle f ^ {\prime} (t) | f (t) \rangle + \langle f (t) | f ^ {\prime} (t) \rangle \\ & = \langle i \overline {{u}} (f (t)) | f (t) \rangle + \langle f (t) | i \overline {{u}} (f (t)) \rangle = 0 \end{array}
\]

因为 \(\overline{u}\) 是自伴的。因此 f = 0，故 \(\varrho(t)x = \pi(t)x\) 对所有 \(t \in R\)。E 的连续自同态 \(\varrho(t)\) 与 \(\pi(t)\) 在 dom(u) 上一致，因而对所有 \(t \in R\) 相等。于是我们有 \(\pi = \varrho\)；由于 \(\overline{u}\) 是 \(\pi\) 的无穷小生成元（引理 14），我们有 \(\overline{u} = u\)，这就证明了 u 是自伴的。

现在我们可以证明定理 1 了。

映射 \(\sigma\) 与 \(\tau\) 是完全有定义的（分别是引理 15 与引理 11）。

设 \(\varrho\) 是 \(\mathbf{R}\) 的一个酉表示。用 \(u\) 表示它的无穷小生成元。关系式 \(\varrho(t) = e^{itu}\)（对所有 \(t \in \mathbf{R}\)）（引理 15）证明了 \(\tau \circ \sigma\) 是恒同映射。

设 u 是 E 上的一个自伴部分算子。引理 14 表明酉表示 \(t \mapsto e^{itu}\) 的无穷小生成元等于 u，因而 \(\sigma \circ \tau\) 是恒同映射。

设 u 是 E 上的一个自伴部分算子，\(\varrho(t) = e^{itu}\)（\(t \in R\)）是与之相伴的 R 在 E 中的酉表示。

设 \(x \in \text{dom}(u)\)。此时所满足的方程 \(\partial_{t} \varrho(t)x = i u(\varrho(t)x)\)（引理 13, b)）称为薛定谔方程。

推论. --- 设 \(\varrho\) 是 \(\mathbf{R}\) 在希尔伯特空间 E 中的一个酉表示。存在一个局部紧空间 X、X 上的一个正测度 \(\mu\) 以及 X 上的一个实值连续函数 \(g\)，使得 \(\varrho\) 同构于表示 \(\pi\)——\(\mathbf{R}\) 在 \(\mathrm{L}^2 (\mathrm{X},\mu)\) 中的——它由 \(\pi (t)f = e^{itg}f\) 定义（对所有 \(t\in \mathbf{R}\) 与每个 \(f\in \mathrm{L}^2 (\mathrm{X},\mu)\)）。

设 \(u\) 是 E 上使得 \(\varrho(t) = e^{itu}\)（对所有 \(t \in \mathbf{R}\)）的自伴算子（定理 1）。存在一个局部紧空间 X、X 上的一个正测度 \(\mu\)、一个等距同构 \(\theta\)（\(L^2(X, \mu)\) 到 E 上的） 以及 X 上的一个实值连续函数 \(g\)，使得 \(u = \theta \circ m_g \circ \theta^{-1}\)（IV, 第 266 页, 定理 1）。断言由公式 \(e^{itu} = \theta \circ e^{itm_g} \circ \theta^{-1} = \theta \circ m_{e^{itg}} \circ \theta^{-1}\)（IV, 第 269 页, 引理 4）得出。

注. --- 假设 u 是 E 的一个自同态。由 \(\varrho(t) = e^{itu}\) 定义的 R 在 E 中的酉表示此时满足不等式 \(\|\varrho(t) - 1_{\mathrm{E}}\| \leqslant |t||u\|\)，对所有 \(t \in R\)，且映射 \(\varrho\)（从 R 到巴拿赫空间 \(\mathcal{L}(\mathrm{E})\) 中的）是线性微分方程

\[
\frac {1}{i} \frac {d \varrho}{d t} = u \circ \varrho
\]

（参见 FVR, IV, 第 26 页, §6）的唯一解。

例. --- 设 \(\varrho\) 是 \(\mathbf{R}\) 在 \(\mathrm{L}^2 (\mathbf{R})\) 中的正则表示。\(\varrho\) 的无穷小生成元是以 \(\mathcal{D}(\mathbf{R})\) 为定义域、由 \(f\mapsto -if'\) 定义的微分算子的闭包。

\section{正定函数}

在本节中，除非另有说明，所考虑的所有向量空间，以及所有希尔伯特空间与代数，都是在 C 上的。

\subsection{普遍正核}

在本章中，X 是一个分离拓扑空间。

定理 1. --- 设 \(f \in \mathcal{C}(X \times X)\)。下列条件等价：

\begin{enumerate}
\def\labelenumi{(\roman{enumi})}
\tightlist
\item
  对 X 的每个紧子集 Y 及 Y 上的每个正测度 \(\mu\)，\(\mathrm{L}^{2}(\mathrm{Y},\mu)\) 的由核 \(f|(\mathrm{Y}\times\mathrm{Y})\)（III, 第 29 页, 定义 1）定义的自同态是正的，换句话说，我们有
\end{enumerate}

\[
\int_ {\mathrm{Y} \times \mathrm{Y}} \overline {{h (x)}} h (y) f (x, y) d (\mu \otimes \mu) (x, y) \geqslant 0
\]

对所有 \(h \in \mathcal{L}^{2}(\mathrm{Y}, \mu)\) 成立；

\begin{enumerate}
\def\labelenumi{(\roman{enumi})}
\setcounter{enumi}{1}
\tightlist
\item
  对每个整数 \(n \in N\)、X 中每个族 \((x_{i})_{0 \leqslant i \leqslant n}\) 以及复数的任意族 \((t_{i})_{0 \leqslant i \leqslant n}\)，我们有
\end{enumerate}

\[
\sum_ {i = 0} ^ {n} \sum_ {j = 0} ^ {n} \bar {t} _ {i} t _ {j} f (x _ {i}, x _ {j}) \geqslant 0;
\]

\begin{enumerate}
\def\labelenumi{(\roman{enumi})}
\setcounter{enumi}{2}
\item
  存在一个复希尔伯特空间 E 和一个具有完全系像的连续映射 \(g \colon X \to E\)，使得 \(f(x, y) = \langle g(x) | g(y) \rangle\) 对所有 \(x\) 与 \(y\)（\(X\) 中的）成立；
\item
  存在一个复希尔伯特空间 E 和一个连续映射 \(g: X \to E\)，使得对所有 X 中的 x 与 y 有 \(f(x, y) = \langle g(x) | g(y) \rangle\)。
\end{enumerate}

(iii) 蕴含 (iv) 是显然的，而通过考虑离散测度 \(\mu\)，即 \(\{0,\ldots,n\}\) 上的计数测度经映射 \(i\mapsto x_{i}\) 的像，以及满足下式的函数 h，可知 (i) 蕴含 (ii)

\[
h(x) = \sum_{\substack{0\leqslant i\leqslant n\\ x_{i} = x}}t_{i}.
\]

我们来证明 (iv) 蕴含 (i)。设存在一个复希尔伯特空间 E 和一个连续映射 \(g \colon \mathrm{X} \to \mathrm{E}\)，使得 \(f(x, y) = \langle g(x) \mid g(y) \rangle\) 对所有 \((x, y) \in \mathrm{X} \times \mathrm{X}\)。设 Y 是 X 的一个紧子集，\(\mu\) 是 Y 上的一个正测度，且 \(h \in \mathcal{L}^{2}(Y, \mu)\)。我们有

\[
\begin{array}{r l} & {\int_ {\mathrm{Y} \times \mathrm{Y}} \overline {{h (x)}} h (y) f (x, y) d (\mu \otimes \mu) (x, y)} \\ & {\qquad = \int_ {\mathrm{Y} \times \mathrm{Y}} \overline {{h (x)}} h (y) \langle g (x) | g (y) \rangle d (\mu \otimes \mu) (x, y)} \\ & {\qquad = \left\langle \int_ {\mathrm{Y}} h (x) g (x) d \mu (x) \Big | \int_ {\mathrm{Y}} h (y) g (y) d \mu (y) \right\rangle \geqslant 0} \end{array}
\]

（INT, V, 第 97 页, § 8, 第 4 目, 命题 9）。

最后我们来证明 (ii) 蕴含 (iii)。设 \(\widetilde{\mathrm{E}}\) 是 X 上具有限支撑的复测度组成的空间。对 \(\mu_{1}\) 与 \(\mu_{2}\)（\(\widetilde{\mathrm{E}}\) 中的），我们令

\[
\langle \mu_ {1} \mid \mu_ {2} \rangle = \int_ {\mathrm{X} \times \mathrm{X}} f (x, y) (\overline {{\mu}} _ {1} \otimes \mu_ {2}) (x, y).
\]

这样在 \(\widetilde{\mathrm{E}}\) 上定义的半双线性形式是一个正埃尔米特形式。事实上，设 \(\mu \in \widetilde{\mathrm{E}}\)；存在 X 中的有限族 \((x_{i})_{0\leqslant i\leqslant n}\) 及复数 \((t_i)_{0\leqslant i\leqslant n}\)，使得 \(\mu = \sum_{i=0}^{n} t_i \varepsilon_{x_i}\)。于是我们有

\[
\langle \mu | \mu \rangle = \sum_ {i = 0} ^ {n} \sum_ {j = 0} ^ {n} \bar {t} _ {i} t _ {j} f (x _ {i}, x _ {j}) \geqslant 0
\]

（由假设）。定义 \(\widetilde{g}\colon\mathbf{X}\to\widetilde{\mathbf{E}}\) 为 \(\widetilde{g}(x)=\varepsilon_{x}\)。映射 \(\widetilde{g}\) 的像生成 \(\widetilde{\mathbf{E}}\)。此外，对每个 \((x,y)\in\mathbf{X}\times\mathbf{X}\)，一方面有 \(f(x,y)=\langle\widetilde{g}(x)|\widetilde{g}(y)\rangle\)，另一方面有

\[
\| \widetilde {g} (x) - \widetilde {g} (y) \| ^ {2} = f (x, x) + f (y, y) - f (x, y) - f (y, x),
\]

这蕴含 \(\widetilde{g}\) 连续，因为 \(f\) 连续。

设 E 是 \(\widetilde{E}\) 的分离完备希尔伯特空间（EVT, V, 第 8 页, 命题 4 的推论），\(g\colon X \to E\) 是 \(\widetilde{g}\) 与典范映射 \(\widetilde{E} \to E\) 的复合。则映射 g 连续，它的像在 E 中是完全系，且 \(f(x,y) = \langle g(x) | g(y) \rangle\) 对所有 \((x,y) \in X \times X\)。

用来证明 (ii) 蕴含 (iii) 的方法称为盖尔范德–奈马克–塞加尔构造（GNS 构造）。

定义 1. --- 若函数 \(f \in \mathcal{C}(\mathrm{X} \times \mathrm{X})\) 满足定理 1 的等价条件，则称它是 X 上的一个普遍正核。

若 f 是 X 上的一个普遍正核，则满足出处同上条件 (iv) 的偶 \((E,g)\) 称为 f 的一个希尔伯特实现；若条件 (iii) 得到满足，则称它是一个循环希尔伯特实现。

用 \(\mathrm{Noy}_{+}(\mathrm{X})\) 表示 X 上普遍正核的集合。

设 X' 是一个分离拓扑空间，\(h: X \to X'\) 是一个连续映射。映射 \(f \mapsto f \circ (h, h)\)（\(\mathcal{C}(X' \times X')\) 到 \(\mathcal{C}(X \times X)\) 中的）通过过渡到子空间诱导出映射 \(\mathrm{Noy}_{+}(X') \to \mathrm{Noy}_{+}(X)\)。

命题 1. --- 设 f 是 X 上的一个普遍正核。设 \((\mathrm{E}_{1}, g_{1})\) 与 \((\mathrm{E}_{2}, g_{2})\) 是 f 的希尔伯特实现。假设 \((\mathrm{E}_{1}, g_{1})\) 是一个循环希尔伯特实现。则存在唯一的连续线性映射 \(u: E_{1} \to E_{2}\) 使得 \(g_{2} = u \circ g_{1}\)。这个映射是等距的。若 \((\mathrm{E}_{2}, g_{2})\) 也是循环的，则 u 是一个同构。

\(u\) 的唯一性由 \(g_{1}\) 的像在 \(\mathrm{E}_{1}\) 中是完全系这一事实得出。设 \(\mathrm{F} = \mathbf{C}^{(\mathrm{X})}\)，\((e_x)_{x\in \mathrm{X}}\) 是 \(\mathrm{F}\) 的典范基。对 \(j = 1\) 与 \(j = 2\)，用 \(u_{j}\) 表示从 \(\mathrm{F}\) 到 \(\mathrm{E}_{j}\) 中的由 \(u_{j}(e_{x}) = g_{j}(x)\) 确定的线性映射，并用 \(\mathrm{F}_{j}\) 表示它的像；空间 \(\mathrm{F}_{1}\) 在 \(\mathrm{E}_{1}\) 中稠密。设 \(t = \sum t_{x}e_{x}\) 是 \(\mathrm{F}\) 的一个元素。我们有

\[
\begin{array}{c} \| u _ {1} (t) \| ^ {2} = \sum_ {x, y} \bar {t} _ {x} t _ {y} \langle g _ {1} (x)   |   g _ {1} (y) \rangle = \sum_ {x, y} \bar {t} _ {x} t _ {y} f (x, y) \\ = \sum_ {x, y} \bar {t} _ {x} t _ {y} \langle g _ {2} (x)   |   g _ {2} (y) \rangle = \| u _ {2} (t) \| ^ {2}. \end{array}
\]

因此，存在从 \(F_{1}\) 到 \(F_{2}\) 中的线性等距映射 v 使得 \(u_{2}=v\circ u_{1}\)，特别地，\(g_{2}(x)=v(g_{1}(x))\) 对所有 \(x\in X\)。由于 \(g_{1}\) 的像在 \(E_{1}\) 中是完全系，这个映射延拓为从 \(E_{1}\) 到 \(E_{2}\) 中的等距线性映射 u，使得 \(g_{2}=u\circ g_{1}\)。

由 I 第 107 页的引理 8，\(u\) 的像在 \(\mathrm{E}_2\) 中是闭的。若 \((\mathrm{E}_2,g_2)\) 是一个循环希尔伯特实现，则 \(u\) 的像还在 \(\mathrm{E}_2\) 中稠密，因而 \(u\) 是一个同构。

命题 2. --- 集合 Noy \(_{+}\) (X) 是空间 \(\mathcal{C}(\text{X} \times \text{X})\) 中的一个自伴锥；它在乘法下稳定，且当给 \(\mathcal{C}(\text{X} \times \text{X})\) 赋以逐点收敛拓扑时它是闭的。

若 \(f \in \mathrm{Noy}_{+}(\mathrm{X})\)，则 \(tf \in \mathrm{Noy}_{+}(\mathrm{X})\) 对每个实数 \(t \geqslant 0\)，且 \(\overline{f} \in \mathrm{Noy}_{+}(\mathrm{X})\)，这是初等的。

若 \((\mathrm{E}_{1}, g_{1})\) 与 \((\mathrm{E}_{2}, g_{2})\) 分别是 X 上普遍正核 \(f_{1}\) 与 \(f_{2}\) 的希尔伯特实现，则偶 \((\mathrm{E}_{1} \oplus \mathrm{E}_{2}, g_{1} + g_{2})\)（相应地，偶 \((\mathrm{E}_{1} \widehat{\otimes}_{2} \mathrm{E}_{2}, g_{1} \otimes g_{2})\)）是 \(f_{1} + f_{2}\)（相应地 \(f_{1}f_{2}\)）的一个希尔伯特实现；因而它们是普遍正核。

\(\mathrm{Noy}_{+}(\mathrm{X})\) 的刻画 (ii)（定理 1）蕴含这个集合在赋以逐点收敛拓扑的 \(\mathcal{C}(\mathrm{X}\times\mathrm{X})\) 中是闭的。

\subsection{关于全纯函数演算的补充注记}

对 C 的每个子集 X，用 \(X^{*}\) 表示 X 在复共轭下的像。设 U 是 C 的一个开子集，\(g: U \to C\) 是一个全纯函数。于是函数 \(f^{*}: z \mapsto \overline{g(\overline{z})}\) 在 \(U^{*}\) 上有定义且全纯。映射 \(f \mapsto f^{*}\) 是从 \(\mathcal{O}(U)\) 到 \(\mathcal{O}(U^{*})\) 上的一个连续双射，使得 \((f_{1}f_{2})^{*} = f_{1}^{*}f_{2}^{*}\) 与 \((f_{1} + f_{2})^{*} = f_{1}^{*} + f_{2}^{*}\) 对 \(f_{1}\) 与 \(f_{2}\)（\(\mathcal{O}(U)\) 中的）成立。

设 C 是 C 的一个紧子集。映射 \(f \mapsto f^{*}\)（\(\mathcal{O}(U)\) 到 \(\mathcal{O}(U^{*})\) 中的）当 U 取遍 C 的包含 C 的开子集时，诱导出空间 \(\mathcal{O}(C)\) 到 \(\mathcal{O}(C^{*})\) 上的一个连续双射（I, 第 49 页, 第 1 目），它也记作 \(f \mapsto f^{*}\)，且满足 \((f_{1}f_{2})^{*} = f_{1}^{*}f_{2}^{*}\) 与 \((f_{1} + f_{2})^{*} = f_{1}^{*} + f_{2}^{*}\) 对 \(f_{1}\) 与 \(f_{2}\)（\(\mathcal{O}(C)\) 中的）。

命题 3. --- 设 A 是一个含单位元的对合巴拿赫代数。设 \(a \in \mathrm{A}\)。\(a^{*}\) 的谱是 \(\mathrm{Sp}_{\mathrm{A}}(a)^{*}\)，即 \(\mathrm{Sp}_{\mathrm{A}}(a)\) 在复共轭下的像。对每个 \(f \in \mathcal{O}(\mathrm{Sp}_{\mathrm{A}}(a))\)，我们有 \(f(a)^{*} = f^{*}(a^{*})\)。

第一个断言由 I 第 97 页得出。由前所述，映射 \(\varphi\)（\(\mathcal{O}(\mathrm{Sp}_{\mathrm{A}}(a))\) 到 A 中的）由 \(f \mapsto (f^{*}(a^{*}))^{*}\) 定义，它是 \(\mathcal{O}(\mathrm{Sp}_{\mathrm{A}}(a))\) 到 A 中的一个连续幺态射，使得在 \(\mathrm{Sp}_{\mathrm{A}}(a)\) 的一个邻域中 \(\mathbf{C}\) 的恒同函数的芽的像等于 \(a\)。因此，\(\varphi\) 就是全纯函数演算的映射 \(f \mapsto f(a)\)（I, 第 74 页, 定理 5）。

引理 1. --- 设 D 是 C 中的单位开圆盘。存在唯一的在 D 上定义的全纯函数 f，使得 \(f(z)^{2} = 1 - z\) 对所有 \(z \in D\)，且 \(f(0) = 1\)。我们有 \(f^{*} = f\)。

幂级数

\[
\sum_ {n = 0} ^ {+ \infty} \binom {1 / 2} {n} (- z) ^ {n}
\]

的收敛半径是 1，其和 \(f\) 定义了 D 上的一个全纯函数（VAR, R1, 第 27 页, 3.2.9），它在 0 处取值 1。它满足 \(f(x)=\sqrt{1-x}\) 对所有 \(x \in D \cap R\)（FVR, III, 第 19 页），因而 \(f(z)^{2} = 1 - z\) 对 \(z \in D\)，因为差 \(f(z)^{2} - (1 - z)\) 是一个各阶导数都在 0 处为零的全纯函数（VAR, R1, 第 27 页, 3.2.5）。按定义即可验证 \(f^{*} = f\)。

设 \(g\) 是 D 上的一个全纯函数，使得 \(g(z)^2 = 1 - z\) 对所有 \(z \in \mathrm{D}\)，且 \(g(0) = 1\)。函数 \(g\) 不取零值，于是 D 上的连续函数 \(f / g\) 取值于 \(\{-1, 1\}\)；由于 D 是连通的且 \(f(0) = g(0)\)，我们有 \(f = g\)。